% संपूर्ण मराठी ग्रंथातील प्रकरण 9: क्रमवर्ती कलन
% हा संपादनयोग्य स्रोतखंड आहे. संकलनासाठी संपूर्ण स्रोतसंग्रहातील
% अक्षररूपे, पूर्वभाग, चिन्हव्याख्या आणि आकृती-स्रोत आवश्यक आहेत.
% मूळ: Open Logic Project; मजकूर आणि रूपांतर CC BY 4.0.
% यंत्रानुवाद, दुरुस्ती आणि तपासणी: OpenAI Codex —
% GPT-5.6 Sol आणि GPT-6 Sol, दोन्ही Ultra effort.
% दुरुस्ती-पुनरावलोकन: GPT-6.1 Sol, Ultra effort.
\chapter{क्रमवर्ती कलन}\label{fol:seq::chap}
\begin{editorial}
  या प्रकरणात अभिजात प्रथम-क्रम तर्कशास्त्रासाठी गेंटझेन यांचे मानक
  क्रमवर्ती कलन LK मांडले आहे. यात आणखी उदाहरणे आणि सराव समाविष्ट करणे
  उपयुक्त ठरेल. सिद्धता-पद्धती म्हणून क्रमवर्ती कलनाशी संबंधित मजकूर
  समाविष्ट करण्यासाठी किंवा वगळण्यासाठी ``prfLK'' हा टॅग वापरा.
\end{editorial}






\section{नियम आणि निष्पत्ती}\label{fol:seq:rul:sec}

पुढील मजकुरात, $\Gamma, \Delta, \Pi, \Lambda$ या वाक्यांच्या सांत
क्रमिका दर्शवतील, असे मानू.

\begin{defn}[क्रमवर्ती]
\emph{क्रमवर्ती} म्हणजे पुढील रूपाची अभिव्यक्ती होय:
\[\Gamma \Sequent \Delta
\]
येथे $\Gamma$ आणि $\Delta$ या $\Lang L$ भाषेतील वाक्यांच्या
सांत (संभवतः रिक्त) क्रमिका आहेत. $\Gamma$ ला \emph{पूर्वांग}, तर
$\Delta$ ला \emph{उत्तरांग} म्हणतात.
\end{defn}

\begin{explain}
क्रमवर्तीमागील अंतःप्रेरित कल्पना अशी आहे: पूर्वांगातील सर्व वाक्ये
सत्य असतील, तर उत्तरांगातील वाक्यांपैकी किमान एक सत्य असते. म्हणजे,
$\Gamma = \tuple{\varphi_1, \dots, \varphi_m}$ आणि
$\Delta = \tuple{\psi_1, \dots, \psi_n}$ असतील, तर $\Gamma \Sequent \Delta$
तेव्हा आणि केवळ तेव्हाच सत्य असते, जेव्हा
\[(\varphi_1 \land \cdots \land \varphi_m) \lif (\psi_1 \lor \cdots \lor
\psi_n)
\]
सत्य असते. याची दोन विशेष प्रकरणे आहेत: $\Gamma$ रिक्त असण्याचे आणि
$\Delta$~रिक्त असण्याचे. $\Gamma$ रिक्त असेल, म्हणजेच $m = 0$ असेल, तर
$\quad \Sequent \Delta$ तेव्हा आणि केवळ तेव्हाच सत्य असते, जेव्हा
$\psi_1 \lor \dots \lor \psi_n$ सत्य असते. $\Delta$ रिक्त असेल, म्हणजेच
$n = 0$ असेल, तर $\Gamma \Sequent \quad$ तेव्हा आणि केवळ तेव्हाच सत्य
असते, जेव्हा $\lnot(\varphi_1 \land \dots \land \varphi_m)$ सत्य असते. एखाद्या
क्रमवर्तीला तेव्हा आणि केवळ तेव्हाच वैध म्हणतो, जेव्हा तिच्याशी संबंधित
वाक्य वैध असते.
\end{explain}

$\Gamma$ ही वाक्यांची क्रमिका असेल, तर $\Gamma$ च्या उजव्या टोकाला
$\varphi$ जोडून मिळणाऱ्या क्रमिकेसाठी आपण $\Gamma, \varphi$ लिहितो (आणि $\Gamma$
च्या डाव्या टोकाला $\varphi$ जोडून मिळणाऱ्या क्रमिकेसाठी $\varphi, \Gamma$
लिहितो). $\Delta$ हीदेखील वाक्यांची क्रमिका असेल, तर $\Gamma,
\Delta$ ही त्या दोन क्रमिकांची जोडणी असते.

\begin{defn}[आरंभीची क्रमवर्ती]
\emph{आरंभीची क्रमवर्ती} म्हणजे, भाषेतील कोणत्याही वाक्य~$\varphi$ साठी,
पुढीलपैकी एका रूपाची क्रमवर्ती होय:
  \begin{enumerate}
    \item $\varphi \Sequent \varphi$
    \item $\quad \Sequent \ltrue$
    \item $\lfalse \Sequent \quad$
  \end{enumerate}.
\end{defn}

क्रमवर्ती कलनातील निष्पत्त्या हे क्रमवर्तींचे विशिष्ट वृक्ष असतात.
त्यांतील सर्वांत वरच्या क्रमवर्ती आरंभीच्या क्रमवर्ती असतात; आणि एखादी
क्रमवर्ती इतर एक किंवा दोन क्रमवर्तींच्या खाली असेल, तर ती निगमनाच्या
एखाद्या नियमानुसार त्यांच्यापासून योग्य रीतीने निष्पन्न झाली पाहिजे.
$\Log{LK}$ चे नियम दोन मुख्य प्रकारांत विभागले आहेत: \emph{तार्किक} नियम
आणि \emph{संरचनात्मक} नियम. खालच्या क्रमवर्तीतील $\varphi$ आणि/किंवा $\psi$
असलेल्या वाक्याच्या मुख्य संकारक वरून तार्किक नियमांना नावे दिली
जातात. प्रत्येक नियमाच्या दोन आवृत्त्या असतात: संकारक असलेले
वाक्य डावीकडे असणारी क्रमवर्ती निष्पन्न करण्यासाठी एक, आणि ते
वाक्य उजवीकडे असणारी क्रमवर्ती निष्पन्न करण्यासाठी दुसरी.









\section{विधानीय नियम}\label{fol:seq:prl:sec}

\subsection{$\lnot$ साठीचे नियम}

\begin{defish}
\Axiom$ \Gamma \fCenter \Delta, \varphi $
\RightLabel{\LeftR{\lnot}}
\UnaryInf$ \lnot \varphi, \Gamma \fCenter \Delta$
\DisplayProof
\hfill
\Axiom$\varphi, \Gamma \fCenter \Delta$
\RightLabel{\RightR{\lnot}}
\UnaryInf$ \Gamma \fCenter \Delta, \lnot \varphi $
\DisplayProof
\end{defish}

\subsection{$\land$ साठीचे नियम}

\begin{defish}\noindent
\begin{tabular}{l}
\Axiom$ \varphi, \Gamma \fCenter \Delta$
\RightLabel{\LeftR{\land}}
\UnaryInf$ \varphi \land \psi, \Gamma \fCenter \Delta$
\DisplayProof
\\[3ex]
\Axiom$\psi, \Gamma \fCenter \Delta$
\RightLabel{\LeftR{\land}}
\UnaryInf$\varphi \land \psi, \Gamma \fCenter \Delta$
\DisplayProof
\end{tabular}
\hfill
\Axiom$\Gamma \fCenter \Delta, \varphi$
\Axiom$ \Gamma \fCenter \Delta, \psi$
\RightLabel{\RightR{\land}}
\BinaryInf$ \Gamma \fCenter \Delta, \varphi \land \psi $
\DisplayProof
\end{defish}
\subsection{$\lor$ साठीचे नियम}
\begin{defish}
\Axiom$\varphi, \Gamma \fCenter \Delta$
\Axiom$\psi, \Gamma \fCenter \Delta$
\RightLabel{\LeftR{\lor}}
\BinaryInf$\varphi \lor \psi, \Gamma \fCenter \Delta$
\DisplayProof
\hfill
\begin{tabular}{r}
\Axiom$\Gamma \fCenter \Delta, \varphi$
\RightLabel{\RightR{\lor}}
\UnaryInf$ \Gamma \fCenter \Delta, \varphi \lor \psi$
\DisplayProof
\\[3ex]
\Axiom$ \Gamma \fCenter \Delta, \psi$
\RightLabel{\RightR{\lor}}
\UnaryInf$ \Gamma \fCenter \Delta, \varphi \lor \psi$
\DisplayProof
\end{tabular}
\end{defish}
\subsection{$\lif$ साठीचे नियम}
\begin{defish}
\Axiom$ \Gamma \fCenter \Delta, \varphi$
\Axiom$ \psi, \Pi \fCenter \Lambda$
\RightLabel{\LeftR{\lif}}
\BinaryInf$ \varphi \lif \psi, \Gamma, \Pi \fCenter \Delta, \Lambda$
\DisplayProof
\hfill
\Axiom$ \varphi, \Gamma \fCenter \Delta, \psi$
\RightLabel{\RightR{\lif}}
\UnaryInf$ \Gamma \fCenter \Delta, \varphi \lif \psi $
\DisplayProof
\end{defish}
\section{संख्यापकांचे नियम}\label{fol:seq:qrl:sec}
\subsection{$\lforall$ साठीचे नियम}
\begin{defish}
\Axiom$ \varphi(t), \Gamma \fCenter \Delta$
\RightLabel{\LeftR{\lforall}}
\UnaryInf$ \lforall[x][\varphi(x)],\Gamma \fCenter \Delta$
\DisplayProof
\hfill
\Axiom$ \Gamma \fCenter \Delta, \varphi(a) $
\RightLabel{\RightR{\lforall}}
\UnaryInf$ \Gamma \fCenter \Delta, \lforall[x][\varphi(x)]$
\DisplayProof
\end{defish}
\LeftR{\lforall} मध्ये, $t$ हे बंद पद आहे (म्हणजे, चरे नसलेले पद).
\RightR{\lforall} मध्ये, $a$~हा असा स्थिरांक आहे की तो
\RightR{\lforall} नियमाच्या खालच्या क्रमवर्तीमध्ये कोठेही येता कामा नये.
\RightR{\forall} अनुमानाचा $a$~हा \emph{आयगेन चर} आहे, असे आपण
म्हणतो.\footnote{वरील नियमात $a$ हा स्थिरांक असला तरी आपण
``आयगेन चर'' ही संज्ञा वापरतो. यामागे ऐतिहासिक कारणे आहेत.}
\enlargethispage{1pt}
\subsection{$\lexists$ साठीचे नियम}
\begin{defish}
\Axiom$ \varphi(a), \Gamma \fCenter \Delta $
\RightLabel{\LeftR{\lexists}}
\UnaryInf$ \lexists[x][\varphi(x)], \Gamma \fCenter \Delta$
\DisplayProof
\hfill
\Axiom$ \Gamma \fCenter \Delta, \varphi(t) $
\RightLabel{\RightR{\lexists}}
\UnaryInf$ \Gamma \fCenter \Delta, \lexists[x][\varphi(x)]$
\DisplayProof
\end{defish}
पुन्हा, $t$~हे बंद पद आहे आणि $a$~हा असा स्थिरांक आहे की तो
\LeftR{\lexists} नियमाच्या खालच्या क्रमवर्तीमध्ये येत नाही.
\LeftR{\lexists} अनुमानाचा $a$~हा \emph{आयगेन चर} आहे, असे आपण म्हणतो.
\RightR{\lforall} किंवा \LeftR{\lexists} अनुमानाच्या खालच्या
क्रमवर्तीमध्ये आयगेन चर येत नसावा, या अटीला \emph{आयगेन चराची अट}
म्हणतात.
\begin{explain}
ही प्रथा आठवा: $\varphi$ हे असे सूत्र असेल की त्यात
चल~$x$ मुक्त असेल, तर हे आपण~$\varphi(x)$ असे लिहून दर्शवतो.
त्याच संदर्भात, मग $\varphi(t)$ हा~$\Subst{\varphi}{t}{x}$ चा संक्षेप असतो.
त्यामुळे \RightR{\lexists} नियम पुढीलप्रमाणेही लिहिता येईल:
\begin{prooftree}
  \Axiom$\Gamma \fCenter \Delta, \Subst{\varphi}{t}{x}$
  \RightLabel{\RightR{\lexists}}
  \UnaryInf$\Gamma \fCenter \Delta, \lexists[x][\varphi]$
\end{prooftree}
लक्षात घ्या की~$\varphi$ मध्ये $t$ आधीच येऊ शकते; उदा., $\varphi$ हे
~$\Atom{\Obj P}{t,x}$ असू शकते. म्हणून,~$\Gamma \Sequent \Delta,
\Atom{\Obj P}{t,t}$ पासून $\Gamma \Sequent \Delta,
\lexists[x][\Atom{\Obj P}{t,x}]$ निष्पन्न करणे हा
\RightR{\lexists} चा योग्य उपयोग आहे---आधारविधानात~$t$ ज्या ठिकाणी
आले आहे, त्यांपैकी एक किंवा अधिक ठिकाणी (म्हणजे, त्या पदाच्या एक किंवा अधिक
आस्थितींच्या जागी) बद्ध चल~$x$ घालता येतो; सर्वच आस्थिती बदलणे आवश्यक नाही. परंतु
\RightR{\lforall} आणि~\LeftR{\lexists} मधील आयगेन चराच्या अटींनुसार
स्थिरांक~$a$ हा~$\varphi$ मध्ये येता कामा नये. म्हणून,
$\RightR{\lforall}$ वापरून $\Gamma \Sequent \Delta, \Atom{\Obj P}{a,a}$
पासून $\Gamma \Sequent \Delta, \lforall[x][\Atom{\Obj P}{a,x}]$
योग्य रीतीने निष्पन्न करता येत नाही.
\end{explain}
\begin{explain}
\RightR{\lexists} आणि \LeftR{\lforall} मध्ये पद~$t$ बंद असण्याव्यतिरिक्त
त्यावर आणखी कोणतेही निर्बंध नाहीत. याउलट, \LeftR{\lexists} आणि \RightR{\lforall}
नियमांमध्ये आयगेन चराच्या अटीनुसार वरच्या क्रमवर्तीमध्ये~$\varphi(a)$ च्या
बाहेर स्थिरांक~$a$ कोठेही येता कामा नये. ही पद्धत निर्दोष आहे,
म्हणजे ती केवळ वैध क्रमवर्तीच निष्पन्न करते, याची खात्री करण्यासाठी ही
अट आवश्यक आहे. ही अट नसती, तर पुढील निष्पत्त्या अनुज्ञात असत्या:
\begin{prooftree}
  \Axiom$\varphi(a) \fCenter \varphi(a)$
  \RightLabel{*\LeftR{\lexists}}
  \UnaryInf$\lexists[x][\varphi(x)] \fCenter \varphi(a)$
  \RightLabel{\RightR{\lforall}}
  \UnaryInf$\lexists[x][\varphi(x)] \fCenter \lforall[x][\varphi(x)]$
  \DisplayProof\bottomAlignProof
  \qquad
  \Axiom$\varphi(a) \fCenter \varphi(a)$
  \RightLabel{*\RightR{\lforall}}
  \UnaryInf$\varphi(a) \fCenter \lforall[x][\varphi(x)]$
  \RightLabel{\LeftR{\lexists}}
  \UnaryInf$\lexists[x][\varphi(x)] \fCenter \lforall[x][\varphi(x)]$
\end{prooftree}
परंतु $\lexists[x][\varphi(x)] \Sequent \lforall[x][\varphi(x)]$ ही क्रमवर्ती
वैध नाही.
\end{explain}
\section{संरचनात्मक नियम}\label{fol:seq:srl:sec}
क्रमवर्तीच्या डाव्या आणि उजव्या बाजूंवरील वाक्यांची मांडणी बदलता येईल,
असे आणखी काही नियम आपल्याला आवश्यक आहेत. तार्किक नियम ज्या वाक्ये वर
लागू होतो, ती आधारविधानात अगदी डाव्या किंवा अगदी उजव्या टोकाला असणे आवश्यक
असते. त्यामुळे वाक्ये योग्य स्थानी हलवता यावीत म्हणून आपल्याला
``अदलाबदल'' नियमाची गरज आहे. तसेच कधी कधी दोन एकसारखी वाक्ये एकात
एकत्र करता येणे आणि दोन्हींपैकी कोणत्याही बाजूला वाक्य जोडता येणे
महत्त्वाचे असते.
\subsection{दुर्बलीकरण}
\begin{defish}
\Axiom$ \Gamma \fCenter \Delta $
\RightLabel{\LeftR{\Weakening}}
\UnaryInf$ \varphi, \Gamma \fCenter \Delta$
\DisplayProof
\hfill
\Axiom$ \Gamma \fCenter \Delta$
\RightLabel{\RightR{\Weakening}}
\UnaryInf$ \Gamma \fCenter \Delta, \varphi$
\DisplayProof
\end{defish}
\subsection{आकुंचन}
\begin{defish}
\Axiom$ \varphi, \varphi, \Gamma \fCenter \Delta $
\RightLabel{\LeftR{\Contraction}}
\UnaryInf$ \varphi, \Gamma \fCenter \Delta$
\DisplayProof
\hfill
\Axiom$ \Gamma \fCenter \Delta, \varphi, \varphi$
\RightLabel{\RightR{\Contraction}}
\UnaryInf$ \Gamma \fCenter \Delta, \varphi$
\DisplayProof
\end{defish}
\subsection{अदलाबदल}
\begin{defish}
\Axiom$ \Gamma, \varphi, \psi, \Pi \fCenter \Delta $
\RightLabel{\LeftR{\Exchange}}
\UnaryInf$ \Gamma, \psi, \varphi, \Pi \fCenter \Delta$
\DisplayProof
\hfill
\Axiom$ \Gamma \fCenter \Delta, \varphi, \psi, \Lambda$
\RightLabel{\RightR{\Exchange}}
\UnaryInf$ \Gamma \fCenter \Delta, \psi, \varphi, \Lambda$
\DisplayProof
\end{defish}
दुर्बलीकरण, आकुंचन आणि अदलाबदल या अनुमानांची मालिका अनेकदा दुहेरी
अनुमानरेषांनी दर्शवली जाते.
पुढील ``कट'' नावाचा नियम काटेकोर अर्थाने आवश्यक नाही; मात्र त्यामुळे
निष्पत्त्या पुन्हा वापरणे आणि एकत्र करणे पुष्कळ सोपे होते.
\begin{defish}
\[
\Axiom$ \Gamma \fCenter \Delta, \varphi$
\Axiom$ \varphi, \Pi \fCenter \Lambda $
\RightLabel{\Cut}
\BinaryInf$ \Gamma, \Pi \fCenter \Delta, \Lambda$
\DisplayProof
\]
\end{defish}









\section{निष्पत्ती}\label{fol:seq:der:sec}

\begin{explain}
आरंभीची क्रमवर्ती कशी असते हे आपण सांगितले आहे आणि निगमन नियम दिले आहेत.
क्रमवर्ती कलनातील निष्पत्त्या यांपासून विगमनाधारित रीतीने निर्माण केल्या
जातात: प्रत्येक निष्पत्ती ही एकतर स्वतःच आरंभीची क्रमवर्ती असते, किंवा
एका अनुमानाच्या वर असलेल्या एक किंवा दोन निष्पत्त्या ने बनलेली असते.
\end{explain}

\begin{defn}[$\Log{LK}$ निष्पत्ती]
क्रमवर्ती~$S$ ची \emph{$\Log{LK}$-निष्पत्ती} म्हणजे पुढील अटी पूर्ण
करणारा क्रमवर्तींचा सांत वृक्ष होय:
\begin{enumerate}
\item वृक्षातील सर्वांत वरच्या क्रमवर्ती आरंभीच्या क्रमवर्ती असतात.
\item वृक्षातील सर्वांत खालची क्रमवर्ती~$S$ असते.
\item वृक्षातील $S$ व्यतिरिक्त प्रत्येक क्रमवर्ती ही निगमन नियमाच्या योग्य
  उपयोजनाचे आधारविधान असते आणि त्या उपयोजनाचा निष्कर्ष वृक्षात त्या
  क्रमवर्तीच्या थेट खाली असतो.
\end{enumerate}
यानंतर $S$ ला त्या निष्पत्तीची \emph{अंतिम क्रमवर्ती} म्हणतो आणि
$S$ ही \emph{$\Log{LK}$ मध्ये निष्पन्न करता येण्याजोगे} आहे (किंवा ती
$\Log{LK}$-निष्पन्न करता येण्याजोगे आहे) असे म्हणतो.
\end{defn}

\begin{ex}
प्रत्येक आरंभीची क्रमवर्ती, उदा., $\chi \Sequent \chi$, ही निष्पत्ती असते.
उदाहरणार्थ, $\LeftR{\Weakening}$ नियम लागू करून आपण तिच्यापासून नवी
निष्पत्ती मिळवू शकतो:
\begin{prooftree}
\Axiom$ \Gamma \fCenter \Delta $
\RightLabel{\LeftR{\Weakening}}
\UnaryInf$ \varphi, \Gamma \fCenter \Delta$
\end{prooftree}
मात्र हा नियम सर्वसाधारण स्वरूपाचा आहे: नियमातील $\varphi$ च्या जागी कोणतेही
वाक्य, उदा., $\theta$ सुद्धा, घेता येते. आधारविधान आपल्या आरंभीच्या
क्रमवर्ती $\chi \Sequent \chi$ शी जुळत असेल, तर $\Gamma$ आणि $\Delta$ ही दोन्ही
केवळ~$\chi$ असतात आणि निष्कर्ष $\theta, \chi \Sequent \chi$ असा असेल. म्हणून पुढील
ही निष्पत्ती आहे:
\begin{prooftree}
\Axiom$ \chi \fCenter \chi $
\RightLabel{\LeftR{\Weakening}}
\UnaryInf$ \theta, \chi \fCenter \chi$
\end{prooftree}
आता आपण दुसरा नियम, समजा $\LeftR{\Exchange}$, लागू करू शकतो; त्यामुळे
डावीकडील दोन वाक्यांची अदलाबदल करता येते. म्हणून पुढीलही योग्य
निष्पत्ती आहे:
\begin{prooftree}
\Axiom$ \chi \fCenter \chi $
\RightLabel{\LeftR{\Weakening}}
\UnaryInf$ \theta, \chi \fCenter \chi$
\RightLabel{\LeftR{\Exchange}}
\UnaryInf$ \chi, \theta \fCenter \chi$
\end{prooftree}
नियमाच्या या उपयोजनात—मूळ नियम पुढीलप्रमाणे दिला होता—
\begin{prooftree}
\Axiom$ \Gamma, \varphi, \psi, \Pi \fCenter \Delta $
\RightLabel{\LeftR{\Exchange}}
\UnaryInf$ \Gamma, \psi, \varphi, \Pi \fCenter \Delta,$
\end{prooftree}
$\Gamma$ आणि $\Pi$ दोन्ही रिक्त ठेवले होते, $\Delta$ हे $\chi$ होते, आणि
$\varphi$ व $\psi$ यांच्या भूमिका अनुक्रमे $\theta$ आणि~$\chi$ यांनी बजावल्या. अगदी
याच प्रकारे पुढीलही निष्पत्ती आहे हे आपल्याला दिसते:
\begin{prooftree}
\Axiom$ \theta \fCenter \theta $
\RightLabel{\LeftR{\Weakening}}
\UnaryInf$ \chi, \theta \fCenter \theta$
\end{prooftree}
आता या दोन निष्पत्त्या घेऊन $\RightR{\land}$ च्या साहाय्याने त्यांना
एकत्र करता येते. तो नियम असा होता:
\begin{prooftree}
\Axiom$\Gamma \fCenter \Delta, \varphi$
\Axiom$ \Gamma \fCenter \Delta, \psi$
\RightLabel{\RightR{\land}}
\BinaryInf$ \Gamma \fCenter \Delta, \varphi \land \psi $
\end{prooftree}
आपल्या उदाहरणात आधारविधाने त्यांवर संपणाऱ्या निष्पत्त्यांच्या अंतिम
क्रमवर्तींशी जुळली पाहिजेत. याचा अर्थ $\Gamma$ हे $\chi, \theta$ आहे, $\Delta$
रिक्त आहे, $\varphi$ हे $\chi$ आहे आणि $\psi$ हे $\theta$ आहे. म्हणून अनुमान योग्य
असेल, तर निष्कर्ष $\chi, \theta \Sequent \chi \land \theta$ हा आहे.
\begin{prooftree}
\Axiom$ \chi \fCenter \chi $
\RightLabel{\LeftR{\Weakening}}
\UnaryInf$ \theta, \chi \fCenter \chi$
\RightLabel{\LeftR{\Exchange}}
\UnaryInf$ \chi, \theta \fCenter \chi$
\Axiom$ \theta \fCenter \theta $
\RightLabel{\LeftR{\Weakening}}
\UnaryInf$ \chi, \theta \fCenter \theta$
\RightLabel{\RightR{\land}}
\BinaryInf$ \chi, \theta \fCenter \chi \land \theta $
\end{prooftree}
अर्थात, आधारविधानांचा क्रम उलटाही करता येतो; मग $\varphi$ हे $\theta$ आणि $\psi$
हे~$\chi$ असेल.
\begin{prooftree}
\Axiom$ \theta \fCenter \theta $
\RightLabel{\LeftR{\Weakening}}
\UnaryInf$ \chi, \theta \fCenter \theta$
\Axiom$ \chi \fCenter \chi $
\RightLabel{\LeftR{\Weakening}}
\UnaryInf$ \theta, \chi \fCenter \chi$
\RightLabel{\LeftR{\Exchange}}
\UnaryInf$ \chi, \theta \fCenter \chi$
\RightLabel{\RightR{\land}}
\BinaryInf$ \chi, \theta \fCenter \theta \land \chi $
\end{prooftree}
\end{ex}









\section{निष्पत्ती: उदाहरणे}\label{fol:seq:pro:sec}

\begin{ex}
क्रमवर्ती $\varphi \land \psi \Sequent \varphi$ ची एक $\Log{LK}$-निष्पत्ती द्या.

हवी असलेली अंतिम क्रमवर्ती निष्पत्तीच्या तळाशी लिहून आपण सुरुवात करतो.
\begin{prooftree}
\AxiomC{}
\UnaryInf$\varphi\land \psi \fCenter \varphi$
\end{prooftree}
यानंतर, या स्वरूपाची खालची क्रमवर्ती कोणत्या प्रकारच्या अनुमानामुळे मिळू शकते
हे शोधावे लागते. ते एखाद्या संरचनात्मक नियमाचे उपयोजन असू शकते; पण आधी
तार्किक नियम शोधणे योग्य ठरते. खालच्या क्रमवर्तीत येणारा एकमेव तार्किक
संयोजक $\land$ आहे; म्हणून आपण $\land$ चा नियम शोधत आहोत. $\land$ हे चिन्ह
पूर्वांगात येत असल्यामुळे आपण \LeftR{\land} नियम पाहतो.
\begin{prooftree}
\AxiomC{}
\RightLabel{\LeftR{\land}}
\UnaryInf$\varphi\land \psi \fCenter \varphi$
\end{prooftree}
\LeftR{\land} अनुमानाची वरची क्रमवर्ती काय असू शकते यासाठी दोन पर्याय आहेत:
वरची क्रमवर्ती $\varphi \Sequent \varphi$ किंवा $\psi \Sequent \varphi$ असू शकते. स्पष्टपणे,
$\varphi \Sequent \varphi$ ही आरंभीची क्रमवर्ती आहे (ही आपल्या दृष्टीने चांगली बाब
आहे), तर $\psi \Sequent \varphi$ ही सर्वसाधारणपणे निष्पन्न करता येत नाही. आपण वरची
क्रमवर्ती भरू:
\begin{prooftree}
\Axiom$\varphi \fCenter \varphi$
\RightLabel{\LeftR{\land}}
\UnaryInf$\varphi\land \psi \fCenter \varphi$
\end{prooftree}
आता आपल्याकडे क्रमवर्ती $\varphi \land \psi \Sequent \varphi$ ची योग्य
$\Log{LK}$-निष्पत्ती आहे.
\end{ex}

\begin{ex}
क्रमवर्ती $\lnot \varphi \lor \psi \Sequent \varphi \lif \psi$ ची एक
$\Log{LK}$-निष्पत्ती द्या.

हवी असलेली अंतिम क्रमवर्ती निष्पत्तीच्या तळाशी लिहून सुरुवात करा.
\begin{prooftree}
\AxiomC{}
\UnaryInf$\lnot \varphi \lor \psi \fCenter \varphi \lif \psi$
\end{prooftree}
ही अंतिम क्रमवर्ती देऊ शकेल असा तार्किक नियम शोधण्यासाठी तिच्यातील तार्किक
संयोजक पाहू: $\lnot$, $\lor$ आणि $\lif$. सध्या आपल्याला फक्त $\lor$ आणि
$\lif$ यांचा विचार करायचा आहे, कारण ती अंतिम क्रमवर्तीतील वाक्यांची
मुख्य संकारक आहेत; $\lnot$ दुसऱ्या संयोजकाच्या व्याप्तीच्या आत आहे,
म्हणून त्याचा विचार आपण नंतर करू. शेवटच्या अनुमानासाठी आपल्याकडे
\LeftR{\lor} नियम आणि \RightR{\lif} नियम हे तार्किक नियमांचे पर्याय आहेत.
खरे तर कोणताही नियम निवडता येईल; पण आपण \RightR{\lif} नियम निवडू (दुसरे
काही कारण नसले तरी त्यामुळे दोन शाखांत विभागणे काही काळ पुढे ढकलता येते).
ही खालची क्रमवर्ती देऊ शकणाऱ्या \RightR{\lif} अनुमानांच्या स्वरूपानुसार ते
पुढीलप्रमाणे असले पाहिजे:
\begin{prooftree}
\AxiomC{}
\UnaryInf$ \varphi, \lnot \varphi \lor \psi \fCenter \psi $
\RightLabel{\RightR{\lif}} \UnaryInf$ \lnot \varphi \lor \psi \fCenter \varphi \lif \psi $
\end{prooftree}

$\lnot \varphi \lor \psi$ हे पूर्वांगाच्या बाहेरील टोकाला हलवले, तर आपण
\LeftR{\lor} नियम लागू करू शकतो. रूपबंधानुसार त्याचे पुढीलप्रमाणे दोन वरच्या
क्रमवर्तींत विभाजन झाले पाहिजे:
\begin{prooftree}
\AxiomC{}
\UnaryInf$\lnot \varphi, \varphi \fCenter \psi$
\AxiomC{}
\UnaryInf$\psi, \varphi \fCenter \psi$
\RightLabel{\LeftR{\lor}}
\BinaryInf$ \lnot \varphi \lor \psi, \varphi \fCenter \psi $
\RightLabel{\LeftR{\Exchange}}
\UnaryInf$ \varphi, \lnot \varphi \lor \psi \fCenter \psi $
\RightLabel{\RightR{\lif}}
\UnaryInf$ \lnot \varphi \lor \psi \fCenter \varphi \lif \psi $
\end{prooftree}
आपण वरच्या दिशेने आरंभीच्या क्रमवर्तींपर्यंत वाट काढण्याचा प्रयत्न करीत आहोत
हे लक्षात ठेवा; आपण बरेच जवळ आलो आहोत असे दिसते!{} उजवी शाखा आरंभीच्या
क्रमवर्तीपासून केवळ एक दुर्बलीकरण आणि एक अदलाबदल दूर आहे; ते केल्यावर ती
पूर्ण होते:
\begin{prooftree}
\AxiomC{}
\UnaryInf$\lnot \varphi, \varphi \fCenter \psi$
\Axiom$\psi \fCenter \psi$
\RightLabel{\LeftR{\Weakening}}
\UnaryInf$\varphi, \psi \fCenter \psi$
\RightLabel{\LeftR{\Exchange}}
\UnaryInf$\psi, \varphi \fCenter \psi$
\RightLabel{\LeftR{\lor}}
\BinaryInf$\lnot \varphi \lor \psi, \varphi \fCenter \psi $
\RightLabel{\LeftR{\Exchange}}
\UnaryInf$ \varphi, \lnot \varphi \lor \psi \fCenter \psi $
\RightLabel{\RightR{\lif}}
\UnaryInf$ \lnot \varphi \lor \psi \fCenter \varphi \lif \psi $
\end{prooftree}

आता डावी शाखा पाहू. कोणत्याही वाक्य मधील एकमेव तार्किक संयोजक हे
पूर्वांगातील वाक्ये मधील $\lnot$ चिन्ह आहे; म्हणून आपण
\LeftR{\lnot} नियमाचे एक उपयोजन शोधत आहोत.
\begin{prooftree}
\AxiomC{}
\UnaryInf$ \varphi \fCenter \psi, \varphi$
\RightLabel{\LeftR{\lnot}}
\UnaryInf$\lnot \varphi, \varphi \fCenter \psi$
\Axiom$\psi \fCenter \psi$
\RightLabel{\LeftR{\Weakening}}
\UnaryInf$\varphi, \psi \fCenter \psi$
\RightLabel{\LeftR{\Exchange}}
\UnaryInf$\psi, \varphi \fCenter \psi$
\RightLabel{\LeftR{\lor}}
\BinaryInf$\lnot \varphi \lor \psi, \varphi \fCenter \psi $
\RightLabel{\LeftR{\Exchange}}
\UnaryInf$ \varphi, \lnot \varphi \lor \psi \fCenter \psi $
\RightLabel{\RightR{\lif}}
\UnaryInf$ \lnot \varphi \lor \psi \fCenter \varphi \lif \psi $
\end{prooftree}
उजवी शाखा जशी पूर्ण केली, त्याचप्रमाणे ही डावी शाखासुद्धा पूर्ण होण्यासाठी
आता केवळ एक दुर्बलीकरण आणि एक अदलाबदल करायची आहे.
\begin{prooftree}
\Axiom$\varphi \fCenter \varphi$
\RightLabel{\RightR{\Weakening}}
\UnaryInf$ \varphi \fCenter \varphi, \psi$
\RightLabel{\RightR{\Exchange}}
\UnaryInf$ \varphi \fCenter \psi, \varphi$
\RightLabel{\LeftR{\lnot}}
\UnaryInf$\lnot \varphi, \varphi \fCenter \psi$
\Axiom$\psi \fCenter \psi$
\RightLabel{\LeftR{\Weakening}}
\UnaryInf$\varphi, \psi \fCenter \psi$
\RightLabel{\LeftR{\Exchange}}
\UnaryInf$\psi, \varphi \fCenter \psi$
\RightLabel{\LeftR{\lor}}
\BinaryInf$\lnot \varphi \lor \psi, \varphi \fCenter \psi $
\RightLabel{\LeftR{\Exchange}}
\UnaryInf$ \varphi, \lnot \varphi \lor \psi \fCenter \psi $
\RightLabel{\RightR{\lif}}
\UnaryInf$ \lnot \varphi \lor \psi \fCenter \varphi \lif \psi $
\end{prooftree}
\end{ex}

\begin{ex}
क्रमवर्ती $\lnot \varphi \lor \lnot \psi \Sequent \lnot (\varphi \land \psi)$ ची एक
$\Log{LK}$-निष्पत्ती द्या.

वरील तंत्रे वापरून, हवी असलेली अंतिम क्रमवर्ती तळाशी लिहून आपण सुरुवात करतो.
\begin{prooftree}
\AxiomC{}
\UnaryInf$ \lnot \varphi \lor \lnot \psi \fCenter \lnot (\varphi \land \psi) $
\end{prooftree}
अंतिम क्रमवर्तीतील वाक्यांची उपलब्ध मुख्य संक्रियके $\lor$ चिन्ह आणि
$\lnot$ चिन्ह ही आहेत. येथे \LeftR{\lor} किंवा \RightR{\lnot} यांपैकी कोणताही
नियम लागू केला तरी चालेल; पण दोन शाखांत विभागणे क्षणभर टाळता यावे म्हणून आपण
\RightR{\lnot} नियमापासून सुरुवात करू:
\begin{prooftree}
\AxiomC{}
\UnaryInf$\varphi \land \psi, \lnot \varphi \lor \lnot \psi \fCenter $
\RightLabel{\RightR{\lnot}}
\UnaryInf$\lnot \varphi \lor \lnot \psi \fCenter \lnot (\varphi \land \psi)$
\end{prooftree}
आता \LeftR{\land} नियम पाहायचा की \LeftR{\lor} नियम, असा पर्याय आहे.
\LeftR{\land} नियम लागू केल्यावर काय होते ते पाहू: $\varphi, \lnot \varphi \lor \lnot \psi
\Sequent \quad$ या क्रमवर्तीपासून किंवा $\psi, \lnot \varphi \lor \lnot \psi \Sequent
\quad$ या क्रमवर्तीपासून सुरुवात करण्याचा पर्याय आपल्याकडे आहे. ही
निष्पत्ती $\varphi$ आणि $\psi$ यांच्या दृष्टीने सममित असल्यामुळे आपण पहिला पर्याय
निवडू:
\begin{prooftree}
\AxiomC{}
\UnaryInf$\varphi, \lnot \varphi \lor \lnot \psi \fCenter $
\RightLabel{\LeftR{\land}}
\UnaryInf$\varphi \land \psi, \lnot \varphi \lor \lnot \psi \fCenter $
\RightLabel{\RightR{\lnot}}
\UnaryInf$\lnot \varphi \lor \lnot \psi \fCenter \lnot (\varphi \land \psi)$
\end{prooftree}
निष्पत्ती पुढे भरत राहिल्यावर आपल्यासमोर एक अडचण येते:
\begin{prooftree}
\Axiom$\varphi \fCenter \varphi$
\RightLabel{\LeftR{\lnot}}
\UnaryInf$ \lnot \varphi, \varphi \fCenter$
\AxiomC{}
\RightLabel{?}
\UnaryInf$\varphi \fCenter \psi$
\RightLabel{\LeftR{\lnot}}
\UnaryInf$ \lnot \psi, \varphi \fCenter$
\RightLabel{\LeftR{\lor}}
\BinaryInf$\lnot \varphi \lor \lnot \psi, \varphi \fCenter $
\RightLabel{\LeftR{\Exchange}}
\UnaryInf$\varphi, \lnot \varphi \lor \lnot \psi \fCenter $
\RightLabel{\LeftR{\land}}
\UnaryInf$\varphi \land \psi, \lnot \varphi \lor \lnot \psi \fCenter $
\RightLabel{\RightR{\lnot}}
\UnaryInf$\lnot \varphi \lor \lnot \psi \fCenter \lnot (\varphi \land \psi)$
\end{prooftree}
उजव्या शाखेच्या वरच्या टोकाचे आणखी लघुकरण करता येत नाही आणि संरचनात्मक
अनुमानांच्या साहाय्याने तिला आरंभीची क्रमवर्ती बनवताही येत नाही; म्हणून हा
योग्य मार्ग नाही. त्यामुळे वर \LeftR{\land} नियम लागू करणे ही स्पष्टपणे चूक
होती. मागील अवस्थेकडे परत जाऊन त्याऐवजी \LeftR{\lor} नियम लागू केल्यावर
आपल्याला पुढील निष्पत्ती मिळते:
\begin{prooftree}
\AxiomC{}
\UnaryInf$\lnot \varphi, \varphi \land \psi \fCenter $

\AxiomC{}
\UnaryInf$\lnot \psi, \varphi \land \psi \fCenter $

\RightLabel{\LeftR{\lor}}
\BinaryInf$\lnot \varphi \lor \lnot \psi, \varphi \land \psi \fCenter $
\RightLabel{\LeftR{\Exchange}}
\UnaryInf$\varphi \land \psi, \lnot \varphi \lor \lnot \psi \fCenter $
\RightLabel{\RightR{\lnot}}
\UnaryInf$\lnot \varphi \lor \lnot \psi \fCenter \lnot (\varphi \land \psi)$
\end{prooftree}
पूर्वी केल्याप्रमाणे प्रत्येक शाखा पूर्ण केल्यावर आपल्याला पुढील निष्पत्ती
मिळते:
\begin{prooftree}
\Axiom$ \varphi \fCenter\varphi$
\RightLabel{\LeftR{\land}}
\UnaryInf$\varphi \land \psi \fCenter \varphi$
\RightLabel{\LeftR{\lnot}}
\UnaryInf$\lnot \varphi, \varphi \land \psi \fCenter $

\Axiom$ \psi \fCenter \psi$
\RightLabel{\LeftR{\land}}
\UnaryInf$\varphi \land \psi \fCenter \psi$
\RightLabel{\LeftR{\lnot}}
\UnaryInf$\lnot \psi, \varphi \land \psi \fCenter $

\RightLabel{\LeftR{\lor}}
\BinaryInf$\lnot \varphi \lor \lnot \psi, \varphi \land \psi \fCenter $
\RightLabel{\LeftR{\Exchange}}
\UnaryInf$\varphi \land \psi, \lnot \varphi \lor \lnot \psi \fCenter $
\RightLabel{\RightR{\lnot}}
\UnaryInf$\lnot \varphi \lor \lnot \psi \fCenter \lnot (\varphi \land \psi)$
\end{prooftree}
(या पायऱ्यांत $\lnot$ नियमांपेक्षा खाली $\land$ नियम लागू केले असते, तरीही
आपल्याला योग्य निष्पत्ती मिळाली असती).
\end{ex}

\begin{ex}
आतापर्यंत आपण आकुंचन नियम वापरलेला नाही; परंतु कधी कधी तो आवश्यक असतो. असे
घडणारे एक उदाहरण पाहू. आपल्याला $\quad \Sequent \varphi \lor \lnot \varphi$ सिद्ध
करायचे आहे असे समजा. $\RightR{\lor}$ उलट्या दिशेने लागू केल्यास पुढील दोन
निष्पत्त्या पैकी एक मिळेल:
\begin{prooftree}
\AxiomC{}
\UnaryInf$ \fCenter \varphi$
\RightLabel{\RightR{\lor}}
\UnaryInf$ \fCenter \varphi \lor \lnot \varphi$
\DisplayProof\qquad\bottomAlignProof
\AxiomC{}
\UnaryInf$\varphi \fCenter $
\RightLabel{\RightR{\lnot}}
\UnaryInf$ \fCenter \lnot \varphi$
\RightLabel{\RightR{\lor}}
\UnaryInf$ \fCenter \varphi \lor \lnot \varphi$
\end{prooftree}
यांपैकी कोणतीही अर्थातच आरंभीच्या क्रमवर्तीवर संपत नाही. यासाठीची युक्ती अशी
की आकुंचन नियम आपल्याला वाक्याच्या दोन प्रती एकात एकत्र करू देतो—आणि
सिद्धता शोधताना, म्हणजे तळापासून वर जाताना, आपण आधारविधानात $\varphi \lor \lnot
\varphi$ ची एक प्रत राखून ठेवू शकतो, उदा.,
\begin{prooftree}
\AxiomC{}
\UnaryInf$ \fCenter \varphi \lor \lnot \varphi, \varphi$
\RightLabel{\RightR{\lor}}
\UnaryInf$ \fCenter \varphi \lor \lnot \varphi, \varphi \lor \lnot \varphi$
\RightLabel{\RightR{\Contraction}}
\UnaryInf$ \fCenter \varphi \lor \lnot \varphi$
\end{prooftree}
आता आपण $\RightR{\lor}$ दुसऱ्यांदा लागू करून~$\lnot \varphi$ सुद्धा मिळवू शकतो;
त्यातून एक पूर्ण निष्पत्ती मिळते.
\begin{prooftree}
\Axiom$\varphi \fCenter \varphi$
\RightLabel{\RightR{\lnot}}
\UnaryInf$\fCenter \varphi, \lnot \varphi$
\RightLabel{\RightR{\lor}}
\UnaryInf$\fCenter \varphi, \varphi \lor \lnot \varphi$
\RightLabel{\RightR{\Exchange}}
\UnaryInf$ \fCenter \varphi \lor \lnot \varphi, \varphi$
\RightLabel{\RightR{\lor}}
\UnaryInf$ \fCenter \varphi \lor \lnot \varphi, \varphi \lor \lnot \varphi$
\RightLabel{\RightR{\Contraction}}
\UnaryInf$ \fCenter \varphi \lor \lnot \varphi$
\end{prooftree}
\end{ex}

\begin{prob}
पुढील क्रमवर्तींसाठी निष्पत्त्या तयार करा:
\begin{enumerate}
\item $\varphi \land (\psi \land \chi) \Sequent (\varphi \land \psi) \land \chi$.
\item $\varphi \lor (\psi \lor \chi) \Sequent (\varphi \lor \psi) \lor \chi$.
\item $\varphi \lif (\psi \lif \chi) \Sequent \psi \lif (\varphi \lif \chi)$.
\item $\varphi \Sequent \lnot\lnot \varphi$.
\end{enumerate}
\end{prob}

\begin{prob}
पुढील क्रमवर्तींसाठी निष्पत्त्या तयार करा:
\begin{enumerate}
\item $(\varphi \lor \psi) \lif \chi \Sequent \varphi \lif \chi$.
\item $(\varphi \lif \chi) \land (\psi \lif \chi) \Sequent (\varphi \lor \psi) \lif \chi$.
\item $\Sequent \lnot(\varphi \land \lnot \varphi)$.
\item $\psi \lif \varphi \Sequent \lnot \varphi \lif \lnot \psi$.
\item $\Sequent (\varphi \lif \lnot \varphi) \lif \lnot \varphi$.
\item $\Sequent \lnot(\varphi \lif \psi) \lif \lnot \psi$.
\item $\varphi \lif \chi \Sequent \lnot (\varphi \land \lnot \chi)$.
\item $\varphi \land \lnot \chi \Sequent \lnot (\varphi \lif \chi)$.
\item $\varphi \lor \psi, \lnot \psi \Sequent \varphi$.
\item $\lnot \varphi \lor \lnot \psi \Sequent \lnot(\varphi \land \psi)$.
\item $\Sequent (\lnot \varphi \land \lnot \psi) \lif\lnot(\varphi \lor \psi)$.
\item $\Sequent \lnot(\varphi \lor \psi) \lif (\lnot \varphi \land \lnot \psi)$.
\end{enumerate}
\end{prob}

\begin{prob}
पुढील क्रमवर्तींसाठी निष्पत्त्या तयार करा:
\begin{enumerate}
\item $\lnot(\varphi \lif \psi) \Sequent \varphi$.
\item $\lnot(\varphi \land \psi) \Sequent \lnot \varphi \lor \lnot \psi$.
\item $\varphi \lif \psi \Sequent \lnot \varphi \lor \psi$.
\item $\Sequent \lnot \lnot \varphi \lif \varphi$.
\item $\varphi \lif \psi, \lnot \varphi \lif \psi \Sequent \psi$.
\item $(\varphi \land \psi) \lif \chi \Sequent (\varphi \lif \chi) \lor (\psi \lif \chi)$.
\item $(\varphi \lif \psi) \lif \varphi \Sequent \varphi$.
\item $\Sequent (\varphi \lif \psi) \lor (\psi \lif \chi)$.
\end{enumerate}
(या सर्वांसाठी $\RightR{\Contraction}$~नियम आवश्यक आहे.)
\end{prob}








\section{संख्यापकांसह निष्पत्ती}\label{fol:seq:prq:sec}

\begin{ex}
क्रमवर्ती $\lexists[x][\lnot \varphi(x)]
\Sequent \lnot \lforall[x][\varphi(x)]$ ची एक $\Log{LK}$-निष्पत्ती द्या.

संख्यापक हाताळताना आयगेन चराच्या अटीचा भंग होणार नाही याची खात्री करावी
लागते; काही वेळा त्यासाठी ठरावीक अनुमाने कोणत्या क्रमाने करायची, यात फेरबदल
करावा लागतो. सामान्यतः आयगेन चराची अट लागू असलेले नियम आधी हाताळण्याचा
प्रयत्न केल्यास मदत होते (पूर्ण झालेल्या सिद्धतेत ते खालच्या बाजूस असतील).
तसेच पुढे काय येईल याचा विचार करून आरंभीची क्रमवर्ती कशी असेल याचा अंदाज
लावणे उपयुक्त ठरते. आपल्या बाबतीत ती $\varphi(a) \Sequent \varphi(a)$ यासारखी असावी
लागेल. याचा अर्थ असा की संख्यापकांचे नियम ``उलट्या दिशेने'' लागू करताना
$\lforall$ नियम आणि $\lexists$ नियम या दोन्हींसाठी एकच पद---ज्याला आपण $a$
म्हणू---निवडावे लागेल. प्रत्येक नियमासाठी वेगवेगळी पदे निवडली, तर शेवटी
$\varphi(a) \Sequent \varphi(b)$ यासारखी क्रमवर्ती मिळेल; ती अर्थातच निष्पन्न करता
येत नाही.

नेहमीप्रमाणे आपण पुढील क्रमवर्ती लिहून सुरुवात करतो:
\begin{prooftree}
\AxiomC{}
\UnaryInf$\lexists[x][\lnot \varphi(x)] \fCenter \lnot \lforall[x][\varphi(x)]$
\end{prooftree}
आपण \LeftR{\exists} नियम किंवा \RightR{\lnot} नियम यांपैकी कोणताही लागू करू
शकतो. \LeftR{\exists} नियमाला आयगेन चराची अट लागू असल्यामुळे तो उशिरा
हाताळण्यापेक्षा लवकर हाताळणे योग्य ठरेल; म्हणून तोच नियम आधी लागू करू.
\begin{prooftree}
\AxiomC{}
\UnaryInf$ \lnot \varphi(a) \fCenter \lnot \lforall[x][\varphi(x)]$
\RightLabel{\LeftR{\lexists}}
\UnaryInf$ \lexists[x][\lnot \varphi(x)] \fCenter \lnot \lforall[x][\varphi(x)]$
\end{prooftree}
\LeftR{\lnot} आणि \RightR{\lnot} नियम उलट्या दिशेने लागू केल्यावर पुढील
निष्पत्ती मिळते:
\begin{prooftree}
\AxiomC{}
\UnaryInf$\lforall[x][\varphi(x)] \fCenter \varphi(a)$
\RightLabel{\LeftR{\lnot}}
\UnaryInf$\lnot \varphi(a), \lforall[x][\varphi(x)] \fCenter $
\RightLabel{\LeftR{\Exchange}}
\UnaryInf$\lforall[x][\varphi(x)], \lnot \varphi(a) \fCenter $
\RightLabel{\RightR{\lnot}}
\UnaryInf$ \lnot \varphi(a) \fCenter \lnot \lforall[x] \varphi(x)$
\RightLabel{\LeftR{\lexists}}
\UnaryInf$ \lexists[x] \lnot \varphi(x) \fCenter \lnot \lforall[x] \varphi(x)$
\end{prooftree}
या टप्प्यावर \LeftR{\forall} नियम लागू करणे हाच आपल्यापुढील एकमेव पर्याय
आहे. या नियमाला आयगेन चराचे बंधन लागू नसल्यामुळे आपला मार्ग मोकळा आहे.
आपल्याला आरंभीची क्रमवर्ती (जिचे स्वरूप $\varphi(a) \Sequent \varphi(a)$ असे आहे)
मिळवण्याचा प्रयत्न करायचा आहे, हे आठवा; म्हणून नियम लागू करताना $\varphi$ मध्ये
घालायचे पद म्हणून $a$ निवडले पाहिजे.
\begin{prooftree}
\Axiom$\varphi(a) \fCenter \varphi(a)$
\RightLabel{\LeftR{\lforall}}
\UnaryInf$\lforall[x][\varphi(x)] \fCenter \varphi(a)$
\RightLabel{\LeftR{\lnot}}
\UnaryInf$\lnot \varphi(a), \lforall[x][\varphi(x)] \fCenter $
\RightLabel{\LeftR{\Exchange}}
\UnaryInf$\lforall[x][\varphi(x)], \lnot \varphi(a) \fCenter $
\RightLabel{\RightR{\lnot}}
\UnaryInf$ \lnot \varphi(a) \fCenter \lnot \lforall[x][\varphi(x)]$
\RightLabel{\LeftR{\lexists}}
\UnaryInf$ \lexists[x][ \lnot \varphi(x)] \fCenter \lnot \lforall[x][\varphi(x)]$
\end{prooftree}
विशेषतः संख्यापक हाताळताना, या टप्प्यावर आयगेन चराच्या अटीचा भंग झालेला नाही
ना, हे पुन्हा तपासणे महत्त्वाचे असते. आयगेन चराची अट लागू असलेला आपण वापरलेला
एकमेव नियम \LeftR{\exists} होता आणि आयगेन चर~$a$ त्याच्या खालच्या क्रमवर्तीत
(म्हणजे अंतिम क्रमवर्तीत) येत नाही; म्हणून ही योग्य निष्पत्ती आहे.
\end{ex}

\begin{prob}
पुढील क्रमवर्तींसाठी निष्पत्त्या तयार करा:
\begin{enumerate}
\item $\Sequent (\lforall[x][\varphi(x)] \land \lforall[y][\psi(y)]) \lif
\lforall[z][(\varphi(z) \land \psi(z))]$.
\item $\Sequent (\lexists[x][\varphi(x)] \lor \lexists[y][\psi(y)]) \lif
\lexists[z][(\varphi(z) \lor \psi(z))]$.
\item $\lforall[x][(\varphi(x) \lif \psi)] \Sequent \lexists[y][\varphi(y)] \lif \psi$.
\item $\lforall[x][\lnot \varphi(x)] \Sequent \lnot\lexists[x][\varphi(x)]$.
\item $\Sequent \lnot\lexists[x][\varphi(x)] \lif \lforall[x][\lnot \varphi(x)]$.
\item $\Sequent \lnot\lexists[x][\lforall[y][((\varphi(x,y) \lif \lnot
\varphi(y,y)) \land (\lnot \varphi(y,y) \lif \varphi(x,y)))]]$.
\end{enumerate}
\end{prob}

\begin{prob}
पुढील क्रमवर्तींसाठी निष्पत्त्या तयार करा:
\begin{enumerate}
\item $\Sequent \lnot\lforall[x][\varphi(x)] \lif \lexists[x][\lnot\varphi(x)]$.
\item $(\lforall[x][\varphi(x)] \lif \psi) \Sequent \lexists[y][(\varphi(y) \lif \psi)]$.
\item $\Sequent \lexists[x][(\varphi(x) \lif \lforall[y][\varphi(y)])]$.
\end{enumerate}
(या सर्वांसाठी $\RightR{\Contraction}$~नियम आवश्यक आहे.)
\end{prob}







\begin{editorial}
  या विभागात नैसर्गिक निगमनासाठी सिद्धतायोग्यता संबंध आणि सुसंगतता यांच्या
  व्याख्या एकत्रित केल्या आहेत.
\end{editorial}



\section{सिद्धता-उपपत्तीय संकल्पना}\label{fol:seq:ptn:sec}

\begin{explain}
जशा आपण अनेक महत्त्वाच्या चिन्हार्थविषयक संकल्पना (वैधता, तार्किक निष्पन्नता,
पूर्ततायोग्यता) परिभाषित केल्या आहेत, तशाच त्यांना अनुरूप
\emph{सिद्धता-उपपत्तीय संकल्पना} आता परिभाषित करू. संरचना मध्ये
वाक्यांची पूर्ति होते की नाही, याचा आधार घेऊन या संकल्पना परिभाषित
केल्या जात नाहीत; त्याऐवजी ठरावीक क्रमवर्तींची निष्पन्नता किंवा
अनिष्पन्नता यांचा आधार घेतला जातो. या दोन प्रकारच्या संकल्पना
एकमेकांशी जुळतात, हा एक महत्त्वाचा शोध होता. त्या जुळतात, हाच
\emph{निर्दोषता प्रमेय} आणि \emph{संपूर्णता प्रमेय} यांचा आशय आहे.
\end{explain}

\begin{defn}[प्रमेये]
वाक्य~$\varphi$ हे \emph{प्रमेय} असते, जर $\Log{LK}$ मध्ये
$\quad \Sequent \varphi$ या क्रमवर्तीची निष्पत्ती असेल. $\varphi$ हे प्रमेय
असेल तर आपण $\Proves \varphi$ लिहितो आणि ते प्रमेय नसेल तर $\Proves/ \varphi$ लिहितो.
\end{defn}

\begin{defn}[निष्पन्नता]
वाक्य~$\varphi$ हे वाक्यांच्या संच~$\Gamma$ \emph{पासून
निष्पन्न करता येण्याजोगे} असते, म्हणजे $\Gamma \Proves \varphi$, तेव्हा आणि केवळ तेव्हाच,
जेव्हा एखादा सांत उपसंच~$\Gamma_0 \subseteq \Gamma$ आणि $\Gamma_0$ मधील
वाक्यांची एखादी क्रमिका $\Gamma_0'$ अशी असते की $\Log{LK}$
$\Gamma_0' \Sequent \varphi$ हे निष्पन्न करते. $\varphi$ हे $\Gamma$ पासून
निष्पन्न करता येण्याजोगे नसेल, तर आपण $\Gamma \Proves/ \varphi$ लिहितो.
\end{defn}

आकुंचन, दुर्बलीकरण आणि अदलाबदल नियमांमुळे $\Gamma_0'$ मधील
वाक्यांचा क्रम आणि त्यांची पुनरावृत्ती कितीदा होते हे महत्त्वाचे नसते:
$\Gamma_0' \Sequent \varphi$ ही क्रमवर्ती निष्पन्न करता येण्याजोगे असेल, तर
$\Gamma_0'$ मध्ये असलेली तीच वाक्ये असणाऱ्या कोणत्याही
$\Gamma_0''$ साठी $\Gamma_0'' \Sequent \varphi$ ही क्रमवर्तीसुद्धा निष्पन्न करता
येते. उदाहरणार्थ, $\Gamma_0 = \{\psi, \chi\}$ असेल, तर
$\Gamma_0' = \tuple{\psi, \psi, \chi}$ आणि $\Gamma_0'' =
\tuple{\chi, \chi, \psi}$ या दोन्ही क्रमिकांमध्ये $\Gamma_0$ मधीलच
वाक्ये आहेत. यांपैकी एक क्रमिका असणारी क्रमवर्ती निष्पन्न करता येण्याजोगे असेल,
तर दुसरी असणारी क्रमवर्तीसुद्धा निष्पन्न करता येते; उदा.:
\begin{prooftree}
  \AxiomC{}
  \Deduce$\psi, \psi, \chi \fCenter \varphi$
  \RightLabel{\LeftR{\Contraction}}
  \UnaryInf$\psi, \chi \fCenter \varphi$
  \RightLabel{\LeftR{\Exchange}}
  \UnaryInf$\chi, \psi \fCenter \varphi$
  \RightLabel{\LeftR{\Weakening}}
  \UnaryInf$\chi, \chi, \psi \fCenter \varphi$
\end{prooftree}
यापुढे $\Gamma_0$ हा वाक्यांचा सांत संच असेल, तर
$\Gamma_0 \Sequent \varphi$ याचा अर्थ असा कोणताही क्रमवर्ती घेऊ की ज्याच्या
पूर्वांगात $\Gamma_0$ मधील वाक्यांची क्रमिका आहे; तसेच आवश्यक असल्यास
आकुंचन, अदलाबदल आणि दुर्बलीकरणे गृहीत धरू.

\begin{defn}[सुसंगतता]
वाक्यांचा संच~$\Gamma$ हा \emph{विसंगत} असतो तेव्हा आणि केवळ तेव्हाच, जेव्हा
एखादा सांत उपसंच~$\Gamma_0 \subseteq \Gamma$ असा असतो की $\Log{LK}$
$\Gamma_0 \Sequent \quad$ हे निष्पन्न करते. $\Gamma$ विसंगत नसेल,
म्हणजे प्रत्येक सांत $\Gamma_0 \subseteq \Gamma$ साठी $\Log{LK}$
$\Gamma_0 \Sequent \quad$ हे निष्पन्न करत नसेल, तर त्याला
\emph{सुसंगत} म्हणतो.
\end{defn}

\begin{prop}[परावर्तिता]
\label{fol:seq:ptn:prop:reflexivity}
$\varphi \in \Gamma$ असेल, तर $\Gamma \Proves \varphi$.
\end{prop}

\begin{proof}
$\varphi \Sequent \varphi$ ही आरंभीची क्रमवर्ती निष्पन्न करता येण्याजोगे आहे आणि $\{\varphi\}
\subseteq \Gamma$.
\end{proof}

\begin{prop}[एकस्वनिकता]
\label{fol:seq:ptn:prop:monotonicity}
$\Gamma \subseteq \Delta$ आणि $\Gamma \Proves \varphi$ असेल, तर $\Delta
\Proves \varphi$.
\end{prop}

\begin{proof}
$\Gamma \Proves \varphi$ असे समजा; म्हणजे, एखादा सांत $\Gamma_0
\subseteq \Gamma$ असा आहे की $\Gamma_0 \Sequent \varphi$ ही क्रमवर्ती
निष्पन्न करता येण्याजोगे आहे. $\Gamma \subseteq \Delta$ असल्यामुळे $\Gamma_0$ हा
$\Delta$ चासुद्धा सांत उपसंच आहे. म्हणून $\Gamma_0 \Sequent \varphi$ ची
निष्पत्ती ही $\Delta \Proves \varphi$ हेसुद्धा दाखवते.
\end{proof}

\begin{prop}[संक्रमकता]
\label{fol:seq:ptn:prop:transitivity}
$\Gamma \Proves \varphi$ आणि $\{\varphi\} \cup \Delta \Proves
\psi$ असेल, तर $\Gamma \cup \Delta \Proves \psi$.
\end{prop}

\begin{proof}
$\Gamma \Proves \varphi$ असेल, तर एखादा सांत $\Gamma_0 \subseteq \Gamma$
आणि $\Gamma_0 \Sequent \varphi$ ची निष्पत्ती~$\pi_0$ असते.
$\{\varphi\} \cup \Delta \Proves \psi$ असेल, तर एखाद्या सांत उपसंचासाठी
$\Delta_0 \subseteq \Delta$, $\varphi, \Delta_0 \Sequent \psi$ ची
निष्पत्ती~$\pi_1$ असते. पुढील निष्पत्ती विचारात घ्या:
\begin{prooftree}
  \AxiomC{}
  \RightLabel{$\pi_0$}
  \Deduce$\Gamma_0 \fCenter \varphi$
  \AxiomC{}
  \RightLabel{$\pi_1$}
  \Deduce$\varphi, \Delta_0 \fCenter \psi$
  \RightLabel{\Cut}
  \BinaryInf$\Gamma_0, \Delta_0 \fCenter \psi$
\end{prooftree}
$\Gamma_0 \cup \Delta_0 \subseteq \Gamma \cup \Delta$ असल्यामुळे यावरून
$\Gamma \cup \Delta \Proves \psi$ हे दिसते.
\end{proof}

विशेषतः, याचा अर्थ असा की $\Gamma \Proves \varphi$ आणि $\varphi
\Proves \psi$ असेल, तर $\Gamma \Proves \psi$. तसेच $\varphi_1,
\dots, \varphi_n \Proves \psi$ असेल आणि प्रत्येक~$i$ साठी $\Gamma \Proves \varphi_i$
असेल, तर $\Gamma \Proves \psi$, हेसुद्धा यावरून मिळते.

\begin{prop}
\label{fol:seq:ptn:prop:incons}
$\Gamma$ विसंगत असतो तेव्हा आणि केवळ तेव्हाच, जेव्हा प्रत्येक
  वाक्य~$\varphi$ साठी $\Gamma \Proves {\varphi}$.
\end{prop}

\begin{proof}
सराव.
\end{proof}

\begin{prob}
\ref{fol:seq:ptn:prop:incons} सिद्ध करा.
\end{prob}

\begin{prop}[संहतता]
  \label{fol:seq:ptn:prop:proves-compact}
  \begin{enumerate}
  \item $\Gamma \Proves \varphi$ असेल, तर एखादा सांत उपसंच $\Gamma_0
    \subseteq \Gamma$ असा आहे की $\Gamma_0 \Proves \varphi$.
  \item $\Gamma$ चा प्रत्येक सांत उपसंच सुसंगत असेल, तर $\Gamma$ सुसंगत
    असतो.
  \end{enumerate}
\end{prop}

\begin{proof}
  \begin{enumerate}
    \item $\Gamma \Proves \varphi$ असेल, तर एखादा सांत उपसंच
      $\Gamma_0 \subseteq \Gamma$ असा आहे की $\Gamma_0
      \Sequent \varphi$ या क्रमवर्तीची निष्पत्ती आहे. परिणामी,
      $\Gamma_0 \Proves \varphi$.
    \item $\Gamma$ विसंगत असेल, तर एखादा सांत
      उपसंच~$\Gamma_0 \subseteq \Gamma$ असा आहे की $\Log{LK}$
      $\Gamma_0 \Sequent \quad$ हे निष्पन्न करते. पण मग $\Gamma_0$ हा
      $\Gamma$ चा विसंगत सांत उपसंच आहे.
  \end{enumerate}
\end{proof}









\section{निष्पन्नता आणि सुसंगतता}\label{fol:seq:prv:sec}

आता आपण निष्पन्नता संबंधाचे अनेक गुणधर्म सिद्ध करू. ते स्वतंत्रपणेही
महत्त्वाचे आहेत, पण संपूर्णता प्रमेयाच्या सिद्धतेत प्रत्येकाची भूमिका असेल.

\begin{prop}\label{fol:seq:prv:prop:provability-contr}
  जर $\Gamma \Proves \varphi$ आणि $\Gamma \cup \{\varphi\}$ विसंगत असेल, तर
  $\Gamma$ विसंगत आहे.
\end{prop}

\begin{proof}
अशी सांत $\Gamma_0$ आणि $\Gamma_1 \subseteq \Gamma$ असतात की
$\Log{LK}$ निष्पन्न $\Gamma_0 \Sequent \varphi$ आणि $\varphi, \Gamma_1
\Sequent \quad$. $\Gamma_0 \Sequent \varphi$ ची $\Log{LK}$-निष्पत्ती
~$\pi_0$ आणि $\Gamma_1, \varphi \Sequent \quad$ ची $\Log{LK}$-निष्पत्ती
~$\pi_1$ अशी नावे देऊ. मग आपण निष्पन्न करू शकतो
\begin{prooftree}
\AxiomC{}
\RightLabel{$\pi_0$}
\Deduce$ \Gamma_0 \fCenter \varphi $
\AxiomC{}
\RightLabel{$\pi_1$}
\Deduce$\varphi, \Gamma_1 \fCenter $
\RightLabel{\Cut}
\BinaryInf$ \Gamma_0,\Gamma_1 \fCenter $
\end{prooftree}

कारण $\Gamma_0 \subseteq \Gamma$ आणि $\Gamma_1 \subseteq \Gamma$,
$\Gamma_0 \cup \Gamma_1 \subseteq \Gamma$; म्हणून $\Gamma$ विसंगत आहे.
\end{proof}

\begin{prop}
\label{fol:seq:prv:prop:prov-incons}
$\Gamma \Proves \varphi$ तेव्हा आणि केवळ तेव्हाच, जेव्हा $\Gamma \cup \{\lnot \varphi\}$
विसंगत असते.
\end{prop}

\begin{proof}
प्रथम $\Gamma \Proves \varphi$ असे गृहीत धरा; म्हणजे $\Gamma \Sequent \varphi$ ची
निष्पत्ती~$\pi_0$ आहे. $\LeftR{\lnot}$ नियम जोडल्यावर आपल्याला
$\lnot \varphi, \Gamma
\Sequent \quad$ ची निष्पत्ती मिळते; म्हणजे $\Gamma \cup \{\lnot \varphi\}$
विसंगत आहे.

जर $\Gamma \cup \{\lnot \varphi\}$ विसंगत असेल, तर $\lnot \varphi, \Gamma
\Sequent \quad$ ची निष्पत्ती~$\pi_1$ असते. खालील $\Gamma \Sequent \varphi$
ची निष्पत्ती आहे:
\begin{prooftree}
  \Axiom$\varphi \fCenter \varphi$
  \RightLabel{\RightR{\lnot}}
  \UnaryInf$\fCenter \varphi, \lnot \varphi$
  \AxiomC{}
  \RightLabel{$\pi_1$}
  \Deduce$\lnot \varphi, \Gamma \fCenter$
  \RightLabel{\Cut}
  \BinaryInf$\Gamma \fCenter \varphi$
\end{prooftree}
\end{proof}

\begin{prob}
$\Gamma \Proves \lnot \varphi$ तेव्हा आणि केवळ तेव्हाच, जेव्हा $\Gamma \cup \{\varphi\}$
विसंगत असते, हे सिद्ध करा.
\end{prob}

\begin{prop}\label{fol:seq:prv:prop:explicit-inc}
  जर $\Gamma \Proves \varphi$ आणि $\lnot \varphi \in \Gamma$ असेल, तर $\Gamma$
  विसंगत आहे.
\end{prop}

\begin{proof}
  समजा $\Gamma \Proves \varphi$ आणि $\lnot \varphi \in \Gamma$. मग $\Gamma_0 \Sequent
  \varphi$ ची निष्पत्ती~$\pi$ असते. $\lnot \varphi, \Gamma_0 \Sequent \quad$
  ही क्रमवर्तीही निष्पन्न करता येण्याजोगे आहे:
  \begin{prooftree}
    \AxiomC{}
    \RightLabel{$\pi$}
    \Deduce$\Gamma_0 \fCenter \varphi$
    \Axiom$\varphi \fCenter \varphi$
    \RightLabel{\LeftR{\lnot}}
    \UnaryInf$\lnot \varphi, \varphi \fCenter$
    \RightLabel{\LeftR{\Exchange}}
    \UnaryInf$\varphi, \lnot \varphi \fCenter$
    \RightLabel{\Cut}
    \BinaryInf$\Gamma_0, \lnot \varphi \fCenter$
  \end{prooftree}
  कारण $\lnot \varphi \in \Gamma$ आणि $\Gamma_0 \subseteq \Gamma$, यावरून
  $\Gamma$ विसंगत असल्याचे दिसते.
\end{proof}

\begin{prop}\label{fol:seq:prv:prop:provability-exhaustive}
  जर $\Gamma \cup \{\varphi\}$ आणि $\Gamma \cup \{\lnot \varphi\}$ दोन्ही
  विसंगत असतील, तर $\Gamma$ विसंगत आहे.
\end{prop}

\begin{proof}
अशी सांत संच $\Gamma_0 \subseteq \Gamma$ आणि $\Gamma_1
\subseteq \Gamma$, तसेच $\Log{LK}$-निष्पत्त्या $\pi_0$ आणि $\pi_1$
$\varphi, \Gamma_0 \Sequent \quad$ आणि $\lnot \varphi, \Gamma_1 \Sequent
\quad$ यांचे अनुक्रमे, असतात. मग आपण निष्पन्न करू शकतो
\begin{prooftree}
\AxiomC{}
\RightLabel{$\pi_0$}
\Deduce$ \varphi, \Gamma_0 \fCenter $
\RightLabel{\RightR{\lnot}}
\UnaryInf$ \Gamma_0 \fCenter \lnot \varphi$
\AxiomC{}
\RightLabel{$\pi_1$}
\Deduce$\lnot \varphi, \Gamma_1 \fCenter  $
\RightLabel{\Cut}
\BinaryInf$ \Gamma_0, \Gamma_1 \fCenter $
\end{prooftree}
कारण $\Gamma_0 \subseteq \Gamma$ आणि $\Gamma_1 \subseteq \Gamma$,
$\Gamma_0 \cup \Gamma_1 \subseteq \Gamma$. म्हणून $\Gamma$ विसंगत आहे.
\end{proof}










\section{निष्पन्नता आणि विधानीय संयोजक}\label{fol:seq:ppr:sec}

\begin{explain}
  क्रमवर्ती कलनातील निष्पन्नता संबंध~$\Proves$ हा विधानीय संयोजकांशी
  संबंधित काही मूलभूत तथ्ये सिद्ध करण्यासाठी पुरेसा सामर्थ्यवान आहे; उदाहरणार्थ,
  $\varphi \land \psi
  \Proves \varphi$ आणि $\varphi, \varphi \lif \psi \Proves \psi$ (विध्यात्मक अनुमान). ही
  तथ्ये संपूर्णता प्रमेयाच्या सिद्धतेसाठी आवश्यक आहेत.
\end{explain}

\begin{prop}\label{fol:seq:ppr:prop:provability-land}
  \begin{enumerate}
  \item \label{fol:seq:ppr:prop:provability-land-left} $\varphi \land \psi \Proves
    \varphi$ आणि $\varphi \land \psi \Proves \psi$ ही दोन्ही तथ्ये सत्य आहेत.
  \item \label{fol:seq:ppr:prop:provability-land-right} $\varphi, \psi \Proves \varphi \land
    \psi$.
  \end{enumerate}
\end{prop}

\begin{proof}
  \begin{enumerate}
  \item $\varphi \land \psi \Sequent \varphi$ आणि $\varphi \land \psi \Sequent
    \psi$ या दोन्ही क्रमवर्ती निष्पन्न करता येण्याजोगे आहेत:
    \begin{prooftree}
      \Axiom$\varphi \fCenter \varphi$
      \RightLabel{\LeftR{\land}}
      \UnaryInf$\varphi \land \psi \fCenter \varphi$
      \DisplayProof\qquad\bottomAlignProof
      \Axiom$\psi \fCenter \psi$
      \RightLabel{\LeftR{\land}}
      \UnaryInf$\varphi \land \psi \fCenter \psi$
    \end{prooftree}
    \item $\varphi, \psi \Sequent \varphi \land \psi$ या क्रमवर्तीची निष्पत्ती येथे आहे:
    \begin{prooftree}
      \Axiom$\varphi \fCenter \varphi$
      \Axiom$\psi \fCenter \psi$
      \RightLabel{\RightR{\land}}
      \BinaryInf$\varphi, \psi \fCenter \varphi \land \psi$
    \end{prooftree}
  \end{enumerate}
\end{proof}

\begin{prop}\label{fol:seq:ppr:prop:provability-lor}
  \begin{enumerate}
  \item $\varphi \lor \psi, \lnot \varphi, \lnot \psi$ विसंगत आहे.
  \item $\varphi \Proves \varphi \lor \psi$ आणि $\psi \Proves \varphi \lor \psi$ ही दोन्ही तथ्ये आहेत.
  \end{enumerate}
\end{prop}

\begin{proof}
  \begin{enumerate}
  \item $\varphi \lor \psi, \lnot \varphi,
    \lnot \psi \Sequent$ या क्रमवर्तीची निष्पत्ती देऊ:
    \begin{prooftree}
      \Axiom$\varphi \fCenter \varphi$
      \RightLabel{\LeftR{\lnot}}
      \UnaryInf$\lnot \varphi, \varphi \fCenter$
      \doubleLine
      \UnaryInf$\varphi, \lnot \varphi, \lnot \psi \fCenter$
      \Axiom$\psi \fCenter \psi$
      \RightLabel{\LeftR{\lnot}}
      \UnaryInf$\lnot \psi, \psi \fCenter$
      \doubleLine
      \UnaryInf$\psi, \lnot \varphi, \lnot \psi \fCenter$
      \RightLabel{\LeftR{\lor}}
      \BinaryInf$ \varphi \lor \psi, \lnot \varphi, \lnot \psi \fCenter $
    \end{prooftree}
    (दुहेरी अनुमान-रेषा दुर्बलीकरण, आकुंचन आणि अदलाबदल यांच्या अनेक अनुमानांना
    दर्शवतात, हे लक्षात ठेवा.)
  \item   $\varphi \Sequent \varphi \lor \psi$ आणि $\psi \Sequent \varphi
    \lor \psi$ या दोन्ही क्रमवर्तींना निष्पत्त्या आहेत:
    \begin{prooftree}
      \Axiom$\varphi \fCenter \varphi$
      \RightLabel{\RightR{\lor}}
      \UnaryInf$\varphi \fCenter \varphi \lor \psi$
      \DisplayProof\qquad\bottomAlignProof
      \Axiom$\psi \fCenter \psi$
      \RightLabel{\RightR{\lor}}
      \UnaryInf$\psi \fCenter \varphi \lor \psi$
    \end{prooftree}
  \end{enumerate}
\end{proof}

\begin{prop}\label{fol:seq:ppr:prop:provability-lif}
  \begin{enumerate}
  \item \label{fol:seq:ppr:prop:provability-lif-left}  $\varphi, \varphi \lif \psi \Proves \psi$.
  \item \label{fol:seq:ppr:prop:provability-lif-right}
    $\lnot \varphi \Proves \varphi \lif \psi$ आणि $\psi \Proves \varphi \lif \psi$ ही दोन्ही तथ्ये आहेत.
  \end{enumerate}
\end{prop}

\begin{proof}
  \begin{enumerate}
    \item $\varphi \lif \psi, \varphi \Sequent \psi$ ही क्रमवर्ती निष्पन्न करता येण्याजोगे आहे:
      \begin{prooftree}
        \Axiom$\varphi \fCenter \varphi$
        \Axiom$\psi \fCenter \psi$
        \RightLabel{\LeftR{\lif}}
        \BinaryInf$\varphi \lif \psi, \varphi  \fCenter \psi$
      \end{prooftree}
    \item $\lnot \varphi \Sequent \varphi \lif \psi$ आणि $\psi
      \Sequent \varphi \lif \psi$ या दोन्ही क्रमवर्ती निष्पन्न करता येण्याजोगे आहेत:
      \begin{prooftree}
        \Axiom$\varphi \fCenter \varphi$
        \RightLabel{\LeftR{\lnot}}
        \UnaryInf$\lnot \varphi, \varphi \fCenter$
        \RightLabel{\LeftR{\Exchange}}
        \UnaryInf$\varphi, \lnot \varphi \fCenter$
        \RightLabel{\RightR{\Weakening}}
        \UnaryInf$\varphi, \lnot \varphi \fCenter \psi$
        \RightLabel{\RightR{\lif}}
        \UnaryInf$\lnot \varphi \fCenter \varphi \lif \psi$
        \DisplayProof\qquad\bottomAlignProof
        \Axiom$\psi \fCenter \psi$
        \RightLabel{\LeftR{\Weakening}}
        \UnaryInf$\varphi, \psi \fCenter \psi$
        \RightLabel{\RightR{\lif}}
        \UnaryInf$\psi \fCenter \varphi \lif \psi$
      \end{prooftree}
  \end{enumerate}
\end{proof}








\section{निष्पन्नता आणि संख्यापक}\label{fol:seq:qpr:sec}

\begin{explain}
  संपूर्णता प्रमेयासाठी या विभागात स्थापित केलेली~$\Proves$ संबंधी तथ्ये
  क्रमवर्ती कलनातील नियमांनी निष्पन्न होतात, हेही आवश्यक आहे.
\end{explain}

\begin{thm}
\label{fol:seq:qpr:thm:strong-generalization} जर $c$ हे $\Gamma$ किंवा $\varphi(x)$ मध्ये
न येणारे स्थिर असेल आणि $\Gamma \Proves \varphi(c)$, तर $\Gamma
\Proves \lforall[x][\varphi(x)]$.
\end{thm}

\begin{proof}
$\pi_0$ ही $\Gamma_0 \Sequent \varphi(c)$ या क्रमवर्तीची $\Log{LK}$-निष्पत्ती
असलेली एखादी सांत $\Gamma_0 \subseteq \Gamma$ घ्या.  $\RightR{\lforall}$
अनुमान जोडल्यावर आपल्याला $\Gamma_0 \Sequent
\lforall[x][\varphi(x)]$ ची निष्पत्ती मिळते, कारण $c$ हे $\Gamma$ किंवा
$\varphi(x)$ मध्ये येत नाही आणि म्हणून आयगेन चराची अट पूर्ण होते.
\end{proof}

\begin{prop}
\label{fol:seq:qpr:prop:provability-quantifiers}
\begin{enumerate}
\item $\varphi(t) \Proves \lexists[x][\varphi(x)]$.
\item $\lforall[x][\varphi(x)] \Proves \varphi(t)$.
\end{enumerate}
\end{prop}

\begin{proof}
\begin{enumerate}
\item $\varphi(t) \Sequent \lexists[x][\varphi(x)]$ ही क्रमवर्ती
  निष्पन्न करता येण्याजोगे आहे:
  \begin{prooftree}
    \Axiom$\varphi(t) \fCenter \varphi(t)$
    \RightLabel{\RightR{\lexists}}
    \UnaryInf$\varphi(t) \fCenter \lexists[x][\varphi(x)]$
    \end{prooftree}
\item $\lforall[x][\varphi(x)] \Sequent \varphi(t)$ ही क्रमवर्ती
  निष्पन्न करता येण्याजोगे आहे:
  \begin{prooftree}
    \Axiom$\varphi(t) \fCenter \varphi(t)$
    \RightLabel{\LeftR{\lforall}}
    \UnaryInf$\lforall[x][\varphi(x)] \fCenter \varphi(t)$
    \end{prooftree}
\end{enumerate}
\end{proof}









\section{निर्दोषता}\label{fol:seq:sou:sec}

\begin{explain}
क्रमवर्ती कलनासारखी निष्पत्ती-पद्धत \emph{निर्दोष} असते, जर ती
प्रत्यक्षात धारण न करणाऱ्या गोष्टी निष्पन्न करू शकत नसेल. त्यामुळे
निर्दोषता हा निष्पत्ती-पद्धतींसाठी हमी दिलेल्या सुरक्षिततेचा एक प्रकार
आहे. कोणता सिद्धता-उपपत्तीय गुणधर्म विचारात आहे यावर अवलंबून, उदाहरणार्थ,
आपल्याला पुढील गोष्टी हव्या असतात:
\begin{enumerate}
\item प्रत्येक निष्पन्न करता येण्याजोगे~$\varphi$ हे वैध असते;
\item एखादे वाक्य इतर काही वाक्यांपासून निष्पन्न करता येण्याजोगे असेल, तर ते
  त्यांचे तार्किक परिणामही असते;
\item वाक्यांचा संच विसंगत असेल, तर तो अपूर्ततायोग्य असतो.
\end{enumerate}
हे निष्पत्ती-पद्धतीचे महत्त्वाचे गुणधर्म आहेत. यांपैकी एखादा गुणधर्म
खरा नसेल, तर निष्पत्ती-पद्धतीत दोष आहे—ती खूप जास्त निष्पत्ती
निष्पन्न करेल. म्हणून निष्पत्ती-पद्धतीची निर्दोषता सिद्ध करणे अत्यंत
महत्त्वाचे आहे.

वरील सर्व सिद्धता-उपपत्तीय गुणधर्म क्रमवर्ती कलनातील ठरावीक क्रमवर्तींच्या
निष्पन्नता द्वारे परिभाषित केलेले असल्याने, (1)--(3) सिद्ध करण्यासाठी
निष्पन्न करता येण्याजोगे क्रमवर्तींच्या चिन्हार्थविषयक गुणधर्मांबद्दल काहीतरी सिद्ध करावे
लागते. प्रथम क्रमवर्ती \emph{वैध} असणे म्हणजे काय ते परिभाषित करू आणि मग
प्रत्येक निष्पन्न करता येण्याजोगे क्रमवर्ती वैध असते हे दाखवू. त्यानंतर (1)--(3) या
परिणामाच्या उपपरिणामांप्रमाणे मिळतील.
\end{explain}

\begin{defn}
संरचना~$\Struct
  M$ क्रमवर्ती
$\Gamma \Sequent \Delta$ \emph{पूर्ण करते} तेव्हा आणि केवळ तेव्हाच, जेव्हा
$\Sat/{M}{\varphi}$ हे एखाद्या $\varphi \in \Gamma$ साठी किंवा
$\Sat{M}{\varphi}$ हे एखाद्या $\varphi \in \Delta$ साठी असते.

क्रमवर्ती \emph{वैध} असते तेव्हा आणि केवळ तेव्हाच, जेव्हा प्रत्येक
संरचना~$\Struct
  M$ ते पूर्ण करते.
\end{defn}

\begin{thm}[निर्दोषता]
\label{fol:seq:sou:thm:sequent-soundness} जर $\Log{LK}$ निष्पन्न $\Theta
\Sequent \Xi$, तर $\Theta \Sequent \Xi$ वैध आहे.
\end{thm}

\begin{proof}
  $\Theta \Sequent \Xi$ ची निष्पत्ती~$\pi$ असू द्या. $\pi$ मधील
अनुमानांची संख्या~$n$ यावर विगमनाने पुढे जाऊ.

अनुमानांची संख्या~$0$ असेल, तर $\pi$ मध्ये फक्त आरंभीची क्रमवर्ती असते.
प्रत्येक आरंभीची क्रमवर्ती $\varphi \Sequent \varphi$ स्पष्टपणे वैध असते, कारण प्रत्येक
$\Struct M$ साठी एकतर
  $\Sat/{M}{\varphi}$ किंवा
$\Sat{M}{\varphi}$.

  अनुमानांची संख्या~0 पेक्षा जास्त असेल, तर सर्वांत खालच्या अनुमानाच्या
प्रकारानुसार प्रकरणे वेगळी करू. विगमन गृहीतकानुसार त्या अनुमानाची आधारविधाने
वैध आहेत असे गृहीत धरू, कारण कोणत्याही आधारविधानाच्या निष्पत्तीमधील
अनुमानांची संख्या~$n$ पेक्षा कमी असते.

प्रथम, एकच आधारविधान असलेली संभाव्य अनुमाने पाहू.
\begin{enumerate}
\item शेवटचे अनुमान दुर्बलीकरण आहे. मग $\Theta \Sequent \Xi$ हे एकतर
  $\varphi, \Gamma \Sequent \Delta$ (शेवटचे अनुमान
  \LeftR{\Weakening} असेल) किंवा $\Gamma \Sequent \Delta, \varphi$ (ते
  \RightR{\Weakening} असेल), आणि निष्पत्ती पुढीलपैकी एका रूपात संपते:
  \begin{prooftree}
    \AxiomC{}
    \Deduce$\Gamma \fCenter \Delta$
    \RightLabel{\LeftR{\Weakening}}
    \UnaryInf$\varphi, \Gamma \fCenter \Delta$
    \DisplayProof\qquad\bottomAlignProof
    \AxiomC{}
    \Deduce$\Gamma \fCenter \Delta$
    \RightLabel{\RightR{\Weakening}}
    \UnaryInf$\Gamma \fCenter \Delta, \varphi$
  \end{prooftree}
  विगमन गृहीतकानुसार $\Gamma \Sequent \Delta$ वैध आहे; म्हणजे प्रत्येक
  संरचना~$\Struct{M}$,
  एकतर एखादे $\chi \in \Gamma$ असे आहे की
  $\Sat/{M}{\chi}$, किंवा एखादे $\chi \in
  \Delta$ असे आहे की $\Sat{M}{\chi}$.

  जर $\Sat/{M}{\chi}$ हे एखाद्या $\chi \in \Gamma$ साठी
  असेल, तर $\Theta = \varphi, \Gamma$ असल्याने $\chi \in \Theta$ देखील आहे आणि
  म्हणून $\Sat/{M}{\chi}$ हे एखाद्या $\chi \in \Theta$ साठी
  आहे. त्याचप्रमाणे, जर $\Sat{M}{\chi}$ हे एखाद्या
  $\chi \in \Delta$ साठी असेल, तर $\chi \in \Xi$ आणि
  $\Sat{M}{\chi}$ हे एखाद्या $\chi \in \Xi$ साठी आहे.
  म्हणून $\Theta \Sequent \Xi$ वैध आहे.
\item शेवटचे अनुमान \LeftR{\lnot} आहे. मग त्याचे आधारविधान
  $\Gamma \Sequent \Delta, \varphi$ आणि निष्कर्ष $\lnot \varphi, \Gamma \Sequent
  \Delta$ असतो; म्हणजे निष्पत्ती पुढील रूपात संपते:
  \begin{prooftree}
    \AxiomC{}
    \Deduce$\Gamma \fCenter \Delta, \varphi$
    \RightLabel{\LeftR{\lnot}}
    \UnaryInf$\lnot \varphi, \Gamma \fCenter \Delta$
  \end{prooftree}
  आणि $\Theta = \lnot \varphi, \Gamma$, तर $\Xi = \Delta$.

  विगमन गृहीतकानुसार $\Gamma \Sequent \Delta, \varphi$ वैध आहे; म्हणजे प्रत्येक
  $\Struct{M}$ साठी एकतर (a) एखाद्या
  $\chi \in \Gamma$ साठी
  $\Sat/{M}{\chi}$, किंवा (b) एखाद्या $\chi \in
  \Delta$ साठी, $\Sat{M}{\chi}$, किंवा (c)
  $\Sat{M}{\varphi}$. $\Theta
  \Sequent \Xi$ देखील वैध आहे हे दाखवायचे आहे. $\Struct{M}$
    ही संरचना घ्या. (a) असल्यास,
  एखादे $\chi \in \Gamma$ असे आहे की
  $\Sat/{M}{\chi}$, आणि $\chi \in \Theta$ मध्ये
  देखील आहे. (b) असल्यास, एखादे $\chi \in \Delta$ असे आहे की
  $\Sat{M}{\chi}$, आणि $\chi \in \Xi$ मध्ये
  देखील आहे. शेवटी, जर $\Sat{M}{\varphi}$ असेल, तर
  $\Sat/{M}{\lnot \varphi}$. कारण $\lnot \varphi \in
  \Theta$ मध्ये असल्याने, एखादे $\chi \in \Theta$ असे आहे की
  $\Sat/{M}{\chi}$. म्हणून $\Theta
  \Sequent \Xi$ वैध आहे.
\item शेवटचे अनुमान \RightR{\lnot} आहे: सरावासाठी.
\item शेवटचे अनुमान \LeftR{\land} आहे. याचे दोन प्रकार आहेत: आधारविधानाच्या
  डाव्या बाजूला असलेल्या $\varphi$ किंवा $\psi$ पासून $\varphi
  \land \psi$ डावीकडे निष्पन्न केले जाऊ शकते. पहिल्या प्रकारात $\pi$ पुढील रूपात
  संपते:
  \begin{prooftree}
    \AxiomC{}
    \Deduce$\varphi, \Gamma \fCenter \Delta$
    \RightLabel{\LeftR{\land}}
    \UnaryInf$\varphi \land \psi, \Gamma \fCenter \Delta$
  \end{prooftree}
  आणि $\Theta = \varphi \land \psi, \Gamma$, तर $\Xi = \Delta$. आता
  संरचना~$\Struct
  M$ घ्या. विगमन गृहीतकानुसार
  $\varphi, \Gamma \Sequent \Delta$ वैध असल्याने, (a)
  $\Sat/{M}{\varphi}$, (b)
  $\Sat/{M}{\chi}$ हे एखाद्या $\chi \in \Gamma$ साठी, किंवा
  (c)~$\Sat{M}{\chi}$ हे एखाद्या $\chi \in \Delta$ साठी.
  (a) प्रकरणात $\Sat/{M}{\varphi \land \psi}$, म्हणून
  $\chi \in \Theta$ (म्हणजे $\varphi \land \psi$) असे आहे की
  $\Sat/{M}{\chi}$. (b) प्रकरणात एखादे $\chi
  \in \Gamma$ असे आहे की $\Sat/{M}{\chi}$ आणि
  $\chi \in \Theta$ देखील आहे. (c) प्रकरणात एखादे $\chi \in \Delta$ असे आहे
  की $\Sat{M}{\chi}$ आणि $\Xi = \Delta$ असल्याने
  $\chi \in \Xi$ देखील आहे. म्हणून प्रत्येक प्रकरणात
  $\Struct
    M$ $\varphi \land \psi, \Gamma \Sequent
  \Delta$ पूर्ण करते. $\Struct{M}$ स्वेच्छ
  असल्याने $\Gamma \Sequent \Delta$ वैध आहे. $\varphi \land \psi$ हे $\psi$ पासून
  निष्पन्न होण्याचे प्रकरणही असेच हाताळता येते; फक्त $\varphi$ ऐवजी $\psi$ घ्या.
\item शेवटचे अनुमान \RightR{\lor} आहे. याचे दोन प्रकार आहेत: आधारविधानाच्या
  उजव्या बाजूला असलेल्या $\varphi$ किंवा $\psi$ पासून $\varphi
  \lor \psi$ उजवीकडे निष्पन्न केले जाऊ शकते. पहिल्या प्रकारात $\pi$ पुढील रूपात
  संपते:
  \begin{prooftree}
    \AxiomC{}
    \Deduce$\Gamma \fCenter \Delta, \varphi$
    \RightLabel{\RightR{\lor}}
    \UnaryInf$\Gamma \fCenter \Delta, \varphi \lor \psi$
  \end{prooftree}
  येथे $\Theta = \Gamma$ आणि $\Xi = \Delta, \varphi \lor \psi$. आता
  संरचना~$\Struct{M}$ घ्या.
  $\Gamma \Sequent \Delta, \varphi$ वैध असल्याने, (a)
  $\Sat{M}{\varphi}$, (b)
  $\Sat/{M}{\chi}$ हे एखाद्या $\chi \in \Gamma$ साठी, किंवा
  (c)~$\Sat{M}{\chi}$ हे एखाद्या $\chi \in \Delta$ साठी.
  (a) प्रकरणात $\Sat{M}{\varphi \lor \psi}$. (b) प्रकरणात
  एखादे $\chi \in \Gamma$ असे आहे की
  $\Sat/{M}{\chi}$. (c) प्रकरणात एखादे $\chi
  \in \Delta$ असे आहे की $\Sat{M}{\chi}$. म्हणून
  प्रत्येक प्रकरणात $\Struct{M}$
  $\Gamma \Sequent \Delta, \varphi \lor \psi$, म्हणजेच $\Theta \Sequent \Xi$ पूर्ण
  करते. $\Struct{M}$ स्वेच्छ असल्याने
  $\Theta \Sequent \Xi$ वैध आहे. $\varphi \lor \psi$ हे $\psi$ पासून निष्पन्न होण्याचे
  प्रकरणही असेच हाताळता येते; फक्त $\varphi$ ऐवजी $\psi$ घ्या.
\item शेवटचे अनुमान \RightR{\lif} आहे. मग $\pi$ पुढील रूपात संपते:
  \begin{prooftree}
    \AxiomC{}
    \Deduce$\varphi, \Gamma \fCenter \Delta, \psi$
    \RightLabel{\RightR{\lif}}
    \UnaryInf$\Gamma \fCenter \Delta, \varphi \lif \psi$
  \end{prooftree}
  पुन्हा, विगमन गृहीतकानुसार आधारविधान वैध आहे; निष्कर्षही वैध आहे हे
  दाखवायचे आहे. $\Struct{M}$ स्वेच्छ घ्या.
  $\varphi, \Gamma \Sequent \Delta, \psi$ वैध असल्याने पुढीलपैकी किमान एक प्रकरण
  लागू होते: (a) $\Sat/{M}{\varphi}$, (b)
  $\Sat{M}{\psi}$, (c)
  $\Sat/{M}{\chi}$ हे एखाद्या~$~\chi \in \Gamma$ साठी, किंवा
  (d) $\Sat{M}{\chi}$ हे एखाद्या $\chi \in \Delta$ साठी.
  (a) आणि (b) प्रकरणांत $\Sat{M}{\varphi \lif \psi}$
  आणि म्हणून $\chi \in \Delta, \varphi \lif \psi$ असे आहे की
  $\Sat{M}{\chi}$. (c) प्रकरणात एखाद्या $\chi \in
  \Gamma$ साठी $\Sat/{M}{\chi}$. (d) प्रकरणात
  एखाद्या $\chi \in \Delta$ साठी $\Sat{M}{\chi}$. प्रत्येक
  प्रकरणात $\Struct{M}$ $\Gamma
  \Sequent \Delta, \varphi \lif \psi$ पूर्ण करते. $\Struct{M}$
  स्वेच्छ असल्याने $\Gamma \Sequent \Delta, \varphi \lif \psi$ वैध आहे.
\item शेवटचे अनुमान \LeftR{\lforall} आहे. मग सूत्र~$\varphi(x)$ आणि
  बंद पद~$t$ असे आहेत की $\pi$ पुढील रूपात संपते:
  \begin{prooftree}
    \AxiomC{}
    \Deduce$\varphi(t), \Gamma \fCenter \Delta$
    \RightLabel{\LeftR{\lforall}}
    \UnaryInf$\lforall[x][\varphi(x)], \Gamma \fCenter \Delta$
  \end{prooftree}
  निष्कर्ष $\lforall[x][\varphi(x)], \Gamma
  \Sequent \Delta$ वैध आहे हे दाखवायचे आहे. संरचना~$\Struct M$ घ्या.
  आधारविधान $\varphi(t), \Gamma \Sequent \Delta$ वैध असल्याने, (a)
  $\Sat/{M}{\varphi(t)}$, (b) $\Sat/{M}{\chi}$ हे एखाद्या $\chi \in \Gamma$ साठी, किंवा
  (c)~$\Sat{M}{\chi}$ हे एखाद्या $\chi \in \Delta$ साठी असते. (a) प्रकरणात,
  \ref{fol:syn:sem:prop:quant-terms} नुसार, जर
  $\Sat{M}{\lforall[x][\varphi(x)]}$ असेल, तर $\Sat{M}{\varphi(t)}$. म्हणून
  $\Sat/{M}{\varphi(t)}$ असल्याने $\Sat/{M}{\lforall[x][\varphi(x)]}$. (b) आणि
  (c) प्रकरणांत $\Struct{M}$ देखील $\lforall[x][\varphi(x)], \Gamma
  \Sequent \Delta$ पूर्ण करते. $\Struct M$ स्वेच्छ असल्याने,
  $\lforall[x][\varphi(x)], \Gamma \Sequent \Delta$ वैध आहे.
\item शेवटचे अनुमान \RightR{\lexists} आहे: सरावासाठी.
\item शेवटचे अनुमान \RightR{\lforall} आहे. मग
  सूत्र~$\varphi(x)$ आणि स्थिरांक~$a$ असे आहेत की $\pi$ पुढील रूपात
  संपते:
  \begin{prooftree}
    \AxiomC{}
    \Deduce$\Gamma \fCenter \Delta, \varphi(a)$
    \RightLabel{\RightR{\lforall}}
    \UnaryInf$\Gamma \fCenter \Delta, \lforall[x][\varphi(x)]$
  \end{prooftree}
  येथे आयगेन चराची अट पूर्ण आहे, म्हणजे $a$ हे $\varphi(x)$, $\Gamma$ किंवा
  $\Delta$ मध्ये येत नाही. विगमन गृहीतकानुसार शेवटच्या अनुमानाचे आधारविधान
  वैध आहे. निष्कर्षही वैध आहे हे दाखवायचे आहे; म्हणजे कोणत्याही
  संरचना~$\Struct M$ साठी (a) $\Sat{M}{\lforall[x][\varphi(x)]}$, (b)
  एखाद्या $\chi \in \Gamma$ साठी $\Sat/{M}{\chi}$, किंवा (c) एखाद्या
  $\chi \in \Delta$ साठी $\Sat{M}{\chi}$.

  $\Struct{M}$ ही स्वेच्छ संरचना घ्या. (b) किंवा (c) लागू असेल तर
  काम झाले; म्हणून दोन्ही लागू नाहीत असे गृहीत धरा: प्रत्येक $\chi \in
  \Gamma$ साठी $\Sat{M}{\chi}$ आणि प्रत्येक $\chi \in \Delta$ साठी
  $\Sat/{M}{\chi}$. आता (a), म्हणजेच $\Sat{M}{\lforall[x][\varphi(x)]}$, हे दाखवायचे
  आहे. \ref{fol:syn:ass:prop:sat-quant} नुसार सर्व चल-विन्यास~$s$ साठी
  $\Sat{M}{\varphi(x)}[s]$ दाखवणे पुरेसे आहे. म्हणून $s$ हा स्वेच्छ चल-विन्यास घ्या.
  $\Struct{M'}$ ही $\Struct{M}$ सारखीच रचना घ्या, फक्त
  $\Assign{a}{M'} = s(x)$ असेल. \ref{fol:syn:ext:cor:extensionality-sent}
  नुसार प्रत्येक $\chi \in \Gamma$ साठी $a$ हे $\Gamma$ मध्ये येत नसल्यामुळे
  $\Sat{M'}{\chi}$, आणि प्रत्येक $\chi \in \Delta$ साठी $\Sat/{M'}{\chi}$.
  पण आधारविधान वैध असल्याने $\Sat{M'}{\varphi(a)}$. \ref{fol:syn:ass:prop:sentence-sat-true}
  नुसार $\varphi(a)$ हे वाक्य असल्यामुळे $\Sat{M'}{\varphi(a)}[s]$.
  आता $\varAssign{s}{s}{x}$ आणि $s(x) = \Value{a}{M'}[s]$, कारण
  $\Struct{M'}$ याच प्रकारे परिभाषित केले आहे. म्हणून
  \ref{fol:syn:ext:prop:ext-formulas} लागू होते आणि $\Sat{M'}{\varphi(x)}[s]$
  मिळते. $a$ हे $\varphi(x)$ मध्ये येत नसल्याने
  \ref{fol:syn:ext:prop:extensionality} नुसार $\Sat{M}{\varphi(x)}[s]$.
  $s$ स्वेच्छ असल्याने सर्व चल-विन्यासांसाठी $\Sat{M}{\varphi(x)}[s]$ सिद्ध झाले.
\item शेवटचे अनुमान \LeftR{\lexists} आहे: सरावासाठी.

\end{enumerate}
आता दोन आधारविधाने असलेली संभाव्य अनुमाने पाहू.
\begin{enumerate}
\item शेवटचे अनुमान कट आहे; म्हणून $\pi$ पुढील रूपात संपते:
  \begin{prooftree}
    \AxiomC{}
    \Deduce$\Gamma \fCenter \Delta, \varphi$
    \AxiomC{}
    \Deduce$\varphi, \Pi \fCenter \Lambda$
    \RightLabel{\Cut}
    \BinaryInf$\Gamma, \Pi \fCenter \Delta, \Lambda$
  \end{prooftree}
  $\Struct{M}$ ही संरचना घ्या. विगमन गृहीतकानुसार आधारविधाने वैध असल्याने
  $\Struct{M}$ दोन्ही आधारविधाने पूर्ण करते.
  दोन प्रकरणे पाहू: (a) $\Sat/{M}{\varphi}$ आणि (b)
  $\Sat{M}{\varphi}$. (a) प्रकरणात डावे आधारविधान पूर्ण
  करण्यासाठी $\Struct{M}$ ने
  $\Gamma \Sequent \Delta$ पूर्ण केले पाहिजे; मग ते निष्कर्षही पूर्ण करते.
  (b) प्रकरणात उजवे आधारविधान पूर्ण करण्यासाठी
  $\Struct{M}$ ने $\Pi \setminus \Lambda$ पूर्ण
  केले पाहिजे. पुन्हा, $\Struct{M}$ निष्कर्ष पूर्ण करते.
\item शेवटचे अनुमान \RightR{\land} आहे. मग $\pi$ पुढील रूपात संपते:
  \begin{prooftree}
    \AxiomC{}
    \Deduce$\Gamma \fCenter \Delta, \varphi$
    \AxiomC{}
    \Deduce$\Gamma \fCenter \Delta, \psi$
    \RightLabel{\RightR{\land}}
    \BinaryInf$\Gamma \fCenter \Delta, \varphi \land \psi$
  \end{prooftree}
  संरचना~$\Struct
    M$ घ्या. $\Struct{M}$
  ने $\Gamma \Sequent \Delta$ पूर्ण केले तर काम झाले. तसे नसल्याचे गृहीत धरा.
  विगमन गृहीतकानुसार $\Gamma \fCenter
  \Delta, \varphi$ वैध असल्याने,
  $\Sat{M}{\varphi}$. त्याचप्रमाणे, कारण $\Gamma
  \Sequent \Delta, \psi$ वैध असल्याने,
  $\Sat{M}{\psi}$. म्हणून
  $\Sat{M}{\varphi \land \psi}$.
\item शेवटचे अनुमान \LeftR{\lor} आहे: सरावासाठी.
\item शेवटचे अनुमान \LeftR{\lif} आहे. मग $\pi$ पुढील रूपात संपते:
  \begin{prooftree}
    \AxiomC{}
    \Deduce$\Gamma \fCenter \Delta, \varphi$
    \AxiomC{}
    \Deduce$\psi, \Pi \fCenter \Lambda$
    \RightLabel{\LeftR{\lif}}
    \BinaryInf$\varphi \lif \psi, \Gamma, \Pi \fCenter \Delta, \Lambda$
  \end{prooftree}
  पुन्हा,
  संरचना~$\Struct{M}$
  $\Struct{M}$ ने $\Gamma, \Pi \Sequent
  \Delta, \Lambda$ पूर्ण केलेले नाही असे गृहीत धरा. आपल्याला
  $\Sat/{M}{\varphi \lif \psi}$ दाखवायचे आहे.
  $\Struct{M}$ ने $\Gamma,
  \Pi \Sequent \Delta, \Lambda$ पूर्ण केलेले नसल्याने ते $\Gamma \Sequent
  \Delta$ किंवा $\Pi \Sequent \Lambda$ यांपैकी कोणतेही पूर्ण करत नाही. $\Gamma
  \Sequent \Delta, \varphi$ वैध असल्याने $\Sat{M}{\varphi}$.
  $\psi, \Pi \Sequent \Lambda$ वैध असल्याने
  $\Sat/{M}{\psi}$. म्हणून
  $\Sat/{M}{\varphi \lif \psi}$, जे दाखवायचे होते तेच.
\end{enumerate}
\end{proof}


\begin{prob}
\ref{fol:seq:sou:thm:sequent-soundness} ची सिद्धता पूर्ण करा.
\end{prob}




\begin{cor}
\label{fol:seq:sou:cor:weak-soundness}
जर $\Proves \varphi$, तर $\varphi$ हे वैध आहे.
\end{cor}

\begin{cor}
\label{fol:seq:sou:cor:entailment-soundness}
जर $\Gamma \Proves \varphi$, तर $\Gamma \Entails \varphi$.
\end{cor}

\begin{proof}
जर $\Gamma \Proves \varphi$ असेल, तर $\Gamma_0 \subseteq \Gamma$ असा काही
मर्यादित संच आणि $\Gamma_0 \Sequent \varphi$ ची निष्पत्ती असते. \ref{fol:seq:sou:thm:sequent-soundness}
नुसार, प्रत्येक
संरचना~$\Struct{M}$
किंवा $\psi \in \Gamma_0$ असत्य करते, किंवा $\varphi$ सत्य करते. म्हणून,
जर $\Sat{M}{\Gamma}$ असेल, तर
$\Sat{M}{\varphi}$ देखील असेल.
\end{proof}

\begin{cor}
\label{fol:seq:sou:cor:consistency-soundness}
जर $\Gamma$ पूर्ण करण्यायोग्य असेल, तर तो सुसंगत आहे.
\end{cor}

\begin{proof}
प्रतिलोम विधान सिद्ध करू. $\Gamma$ सुसंगत नाही असे गृहीत धरा. मग
$\Gamma_0 \subseteq \Gamma$ असा काही मर्यादित संच आणि
$\Gamma_0 \Sequent \quad$ ची निष्पत्ती असते. \ref{fol:seq:sou:thm:sequent-soundness}
नुसार $\Gamma_0 \Sequent \quad$ वैध आहे. दुसऱ्या शब्दांत, प्रत्येक
संरचना~$\Struct{M}$ साठी
$\chi \in \Gamma_0$ असे आहे की
$\Sat/{M}{\chi}$; आणि $\Gamma_0 \subseteq \Gamma$
असल्याने ते $\chi$ $\Gamma$ मध्येही आहे. त्यामुळे कोणतेही
$\Struct{M}$ $\Gamma$ पूर्ण करत नाही, आणि
$\Gamma$ पूर्ण करण्यायोग्य नाही.
\end{proof}









\section{निष्पत्ती सह एकरूपता}\label{fol:seq:ide:sec}

निष्पत्त्या आणि एकरूपता साठी अतिरिक्त आरंभीच्या क्रमवर्ती
आणि अनुमाननियमांची आवश्यकता असते.

\begin{defn}[$\eq$ साठी आरंभीच्या क्रमवर्ती]
जर $t$ हे बंद पद असेल, तर ${} \Sequent \eq[t][t]$ ही आरंभीची क्रमवर्ती असते.
\end{defn}

$\eq$ साठीचे नियम पुढीलप्रमाणे आहेत ($t_1$ आणि $t_2$ ही बंद पदे आहेत):

\begin{defish}
\Axiom$ \eq[t_1][t_2], \Gamma \fCenter \Delta, \varphi(t_1) $
\RightLabel{$\eq$}
\UnaryInf$\eq[t_1][t_2], \Gamma \fCenter \Delta, \varphi(t_2)$
\DisplayProof
\hfill
\Axiom$\eq[t_1][t_2], \Gamma \fCenter \Delta, \varphi(t_2) $
\RightLabel{$\eq$}
\UnaryInf$\eq[t_1][t_2], \Gamma  \fCenter \Delta, \varphi(t_1)$
\DisplayProof
\end{defish}

\begin{ex}
जर $s$ आणि $t$ ही बंद पदे असतील, तर $\eq[s][t], \varphi(s)
\Proves \varphi(t)$:
\begin{prooftree}
\Axiom$ \varphi(s) \fCenter \varphi(s)$
\RightLabel{\LeftR{\Weakening}}
\UnaryInf$\eq[s][t], \varphi(s)  \fCenter \varphi(s)$
\RightLabel{$\eq$}
\UnaryInf$\eq[s][t], \varphi(s)  \fCenter \varphi(t)$
\end{prooftree}
याला एकरूप वस्तूंच्या प्रतिस्थापनाचे तत्त्व किंवा लाइब्निझचा नियम
म्हणून ओळखले जाऊ शकते.

$\Log{LK}$ सिद्ध करते की $\eq$ सममित आणि संक्रमक आहे:
\begin{prooftree}
\Axiom$ \fCenter \eq[t_1][t_1] $
\RightLabel{\LeftR{\Weakening}}
\UnaryInf$ \eq[t_1][t_2] \fCenter \eq[t_1][t_1] $
\RightLabel{$\eq$}
\UnaryInf$ \eq[t_1][t_2] \fCenter \eq[t_2][t_1]$
\DisplayProof\qquad\bottomAlignProof
\Axiom$ \eq[t_1][t_2] \fCenter \eq[t_1][t_2] $
\RightLabel{\LeftR{\Weakening}}
\UnaryInf$\eq[t_2][t_3], \eq[t_1][t_2]  \fCenter \eq[t_1][t_2] $
\RightLabel{$\eq$}
\UnaryInf$\eq[t_2][t_3], \eq[t_1][t_2]  \fCenter \eq[t_1][t_3]$
\RightLabel{\LeftR{\Exchange}}
\UnaryInf$\eq[t_1][t_2], \eq[t_2][t_3]  \fCenter \eq[t_1][t_3]$
\end{prooftree}
डावीकडील निष्पत्तीमध्ये सूत्र~$\eq[x][t_1]$ हे आपले
$\varphi(x)$ आहे. उजवीकडे $\varphi(x)$ हे~$\eq[t_1][x]$ असे घेतो.
\end{ex}

\begin{prob}
पुढील क्रमवर्तींच्या निष्पत्त्या द्या:
\begin{enumerate}
\item $\Sequent \lforall[x][\lforall[y][((x = y \land \varphi(x)) \lif \varphi(y))]]$
\item $\lexists[x][\varphi(x)] \land \lforall[y][\lforall[z][((\varphi(y) \land
    \varphi(z)) \lif y = z)]] \Sequent
\lexists[x][(\varphi(x) \land \lforall[y][(\varphi(y) \lif y = x)])]$
\end{enumerate}
\end{prob}









\section{एकरूपता सह निर्दोषता}\label{fol:seq:sid:sec}

\begin{prop}
$\Log{LK}$ मधील आरंभीच्या क्रमवर्ती आणि एकरूपतेच्या नियमांसहची पद्धत निर्दोष आहे.
\end{prop}

\begin{proof}
${} \Sequent \eq[t][t]$ या रूपातील आरंभीच्या क्रमवर्ती वैध आहेत, कारण
प्रत्येक संरचना~$\Struct M$ साठी $\Sat{M}{\eq[t][t]}$ असते. (लक्षात घ्या की
$t$ हे बंद पद आहे असे आपण गृहीत धरतो, म्हणजे त्यात कोणतेही चल नाही;
म्हणून चल-विन्यास अप्रासंगिक असतात.)

निष्पत्तीमधील शेवटचे अनुमान $=$ आहे असे गृहीत धरा. मग आधारविधान
$\eq[t_1][t_2], \Gamma \Sequent \Delta, \varphi(t_1)$ आणि निष्कर्ष
$\eq[t_1][t_2], \Gamma \Sequent \Delta, \varphi(t_2)$ असतो. संरचना~$\Struct M$
घ्या. निष्कर्ष वैध आहे हे दाखवायचे आहे; म्हणजे $\Sat{M}{\eq[t_1][t_2]}$ आणि
$\Sat{M}{\Gamma}$ असल्यास, एकतर $\Sat{M}{\chi}$ हे एखाद्या $\chi \in \Delta$ साठी
असते किंवा $\Sat{M}{\varphi(t_2)}$ असते.

विगमन गृहीतकानुसार आधारविधान वैध आहे. म्हणजे $\Sat{M}{\eq[t_1][t_2]}$ आणि
$\Sat{M}{\Gamma}$ असल्यास, एकतर (a) एखाद्या $\chi \in \Delta$ साठी
$\Sat{M}{\chi}$ किंवा (b) $\Sat{M}{\varphi(t_1)}$. (a) प्रकरणात काम झाले.
(b) प्रकरणाचा विचार करू. $s$ हा चल-विन्यास घ्या, ज्यासाठी
$s(x) = \Value{t_1}{M}$. \ref{fol:syn:ass:prop:sentence-sat-true} नुसार,
$\Sat{M}{\varphi(t_1)}[s]$. $\varAssign{s}{s}{x}$ असल्याने,
\ref{fol:syn:ext:prop:ext-formulas} नुसार $\Sat{M}{\varphi(x)}[s]$. कारण
$\Sat{M}{\eq[t_1][t_2]}$, आपल्याला $\Value{t_1}{M} = \Value{t_2}{M}$,
आणि म्हणून $s(x) = \Value{t_2}{M}$ मिळते. \ref{fol:syn:ext:prop:ext-formulas}
पुन्हा लागू केल्याने $\Sat{M}{\varphi(t_2)}[s]$ देखील मिळते.
\ref{fol:syn:ass:prop:sentence-sat-true} नुसार,
$\Sat{M}{\varphi(t_2)}$.
\end{proof}



